\[
d ^ {+} (\mu , \gamma , \nu) = \dim \mathrm{E} ^ {+} (\mu , \gamma , \nu), \quad d ^ {-} (\mu , \gamma , \nu) = \dim \mathrm{E} ^ {-} (\mu , \gamma , \nu).
\]

对所有 \(\lambda \in \mathfrak{h}^*\)，置 \(\lambda^{*} = -w_{0}\lambda\)，其中 \(w_{0}\) 是 W 中把 B 变为 \(-B\) 的元素。证明

\[
d ^ {+} (\mu , \gamma , \nu) = d ^ {-} (\mu , - \gamma^ {*}, \nu^ {*}).
\]

\begin{enumerate}
\def\labelenumi{\alph{enumi})}
\setcounter{enumi}{1}
\item
  令 \(\lambda_1, \lambda_2 \in \mathrm{P}_{++}\)，V 为 \(\mathfrak{g}\)-模 \(\mathrm{Hom}_k(\mathrm{E}(\lambda_1^*), \mathrm{E}(\lambda_2))\)，U 为由 \(\varphi \in V\) 中满足 \(Y_\alpha. \varphi = 0\)（对所有 \(\alpha \in B\)）者组成的集合，\(w\) 为 \(\mathrm{E}(\lambda_1^*)\) 中的本原向量。证明 \(\varphi \mapsto \varphi(w)\) 是从 U 到由 \(v \in \mathrm{E}(\lambda_2)\) 中满足 \(Y_\alpha^{\lambda_1^*(H_\alpha)+1}. v = 0\)（对所有 \(\alpha \in B\)）者组成的集合的同构。（为证满射，使用§7的习题15。）
\item
  采用命题2的记号，证明对 \(\lambda_1, \lambda_2, \lambda \in \mathrm{P}_{++}\)，
\end{enumerate}

\[
\begin{array}{r} m (\lambda_ {1}, \lambda_ {2}, \lambda) = d ^ {+} (\lambda , \lambda_ {2} - \lambda_ {1} ^ {*}, \lambda_ {1} ^ {*}) = d ^ {-} (\lambda , \lambda_ {1} - \lambda_ {2} ^ {*}, \lambda_ {1}) \\ = d ^ {+} (\lambda_ {1}, \lambda - \lambda_ {2}, \lambda_ {2}) = d ^ {-} (\lambda_ {1}, \lambda_ {2} ^ {*} - \lambda^ {*}, \lambda_ {2} ^ {*}). \end{array}
\]

（注意 \(m(\lambda_1,\lambda_2,\lambda)\) 是 \(\mathfrak{g}\)-不变元在

\[
\mathrm{E} (\lambda) ^ {*} \otimes \mathrm{E} (\lambda_ {1}) \otimes \mathrm{E} (\lambda_ {2}),
\]

中所成空间的维数，因而等于 \(m(\lambda_1^*, \lambda, \lambda_2)\)。）

\(d)\) 令 \(\lambda_1, \lambda_2 \in \mathrm{P}_{++}\)，并令 \(\lambda\) 为 \(\mathrm{P}_{++} \cap \mathrm{W}.\) ( \(\lambda_1 - \lambda_2^*\) ) 的唯一元素。我们有

\[
m (\lambda_ {1}, \lambda_ {2}, \lambda) = 1.
\]

（使用 c)。）由此推出 \(\mathrm{E}(\lambda_1) \otimes \mathrm{E}(\lambda_2)\) 含有唯一同构于 \(\mathrm{E}(\lambda)\) 的子模。

\begin{enumerate}
\def\labelenumi{\alph{enumi})}
\setcounter{enumi}{4}
\tightlist
\item
  保留 \(d\) ) 的记号。证明下列条件等价：
\end{enumerate}

\begin{enumerate}
\def\labelenumi{(\roman{enumi})}
\item
  \(\lambda = \lambda_{1} + \lambda_{2}\)
\item
  \(\| \lambda \| = \| \lambda_1 + \lambda_2 \|\)
\item
  \(\lambda_{1}\) 与 \(\lambda_2^*\) 正交
\item
  \(\lambda_{1}\) 与 \(\lambda_2^*\) 不相交（习题13）
\item
  \(\lambda_{1}\) 与 \(\lambda_{2}\) 不相交。
\end{enumerate}

由此断定，\(\mathrm{E}(\lambda_{1})\otimes\mathrm{E}(\lambda_{2})\) 不是单模，除非 \(\lambda_{1}\) 与 \(\lambda_{2}\) 不相交（从而给出§7习题18 f) 的又一个证明）。

\begin{enumerate}
\def\labelenumi{\arabic{enumi})}
\setcounter{enumi}{14}
\tightlist
\item
  置 \(\mathrm{N} = \mathrm{Card}(\mathrm{R}_{+})\)，\(c = \prod_{\alpha \in \mathrm{R}_{+}} \langle \rho, H_{\alpha} \rangle\)，并置 \(d(\lambda) = \dim \mathrm{E}(\lambda)\)（对 \(\lambda \in \mathrm{P}_{++}\)）。证明对任意实数 \(s > 0\)，
\end{enumerate}

\[
\sum_ {\lambda \in \mathrm{P} _ {+ +}} d (\lambda) ^ {- s} \leq \frac {1}{c} \left(\sum_ {m = 1} ^ {\infty} m ^ {- s}\right) ^ {\mathrm{N}}.
\]

由此推出 \(\sum_{\lambda \in \mathrm{P}_{+ + }}d(\lambda)^{-s} <   + \infty\)，若 \(s > 1\)

¶ 16) a) 取 g 为 \(F_{4}\) 型，并采用第VI章图版VIII的记号。当 i = 1, 2, 3, 4 时，置

\[
\mathcal {X} _ {i} = \mathrm{P} _ {+ +} \cap (\varpi_ {i} - \mathrm{Q} _ {+}).
\]

\(\mathrm{E}(\varpi_i)\) 的权集合是诸 \(\mathrm{W}\omega\) 的不交并，其中 \(\omega\) 属于 \(\mathcal{X}_i\)（参见§7命题5 (iv)）。我们有：

\[
\mathcal {X} _ {1} = \{0, \varpi_ {1}, \varpi_ {4} \};
\]

\[
\mathcal {X} _ {2} = \{0, \varpi_ {1}, \varpi_ {2}, \varpi_ {3}, \varpi_ {4}, \varpi_ {1} + \varpi_ {4}, 2 \varpi_ {4} \};
\]

\[
\mathcal {X} _ {3} = \{0, \varpi_ {1}, \varpi_ {3}, \varpi_ {4} \};
\]

\[
\mathcal {X} _ {4} = \{0, \varpi_ {4} \}.
\]

\begin{enumerate}
\def\labelenumi{\alph{enumi})}
\setcounter{enumi}{1}
\tightlist
\item
  借助定理2证明：
\end{enumerate}

\(\dim \mathrm{E}(\varpi_1) = 52,\quad \dim \mathrm{E}(\varpi_2) = 1274,\quad \dim \mathrm{E}(\varpi_3) = 273,\quad \dim \mathrm{E}(\varpi_4) = 26.\)

\begin{enumerate}
\def\labelenumi{\alph{enumi})}
\setcounter{enumi}{2}
\tightlist
\item
  利用第V章§3的命题1以及第VI章的图版，证明
\end{enumerate}

\[
\mathrm{Card} (\mathrm{W} \varpi_ {1}) = 2 ^ {7} 3 ^ {2} 2 ^ {- 3} (3!) ^ {- 1} = 24.
\]

类似地计算 \(\mathrm{Card}(\mathrm{W}\varpi_{2}),\ldots,\mathrm{Card}(\mathrm{W}.2\varpi_{4})\)。由此推出 \(\mathrm{E}(\varpi_{2})\) 的权数为 553；由于该数严格小于 \(\dim\mathrm{E}(\varpi_{2})\)，这些权中有一个的重数 \(\geq2\)。

\begin{enumerate}
\def\labelenumi{\alph{enumi})}
\setcounter{enumi}{3}
\item
  对 \(\varpi_{1},\varpi_{3},\varpi_{4}\) 作类似的计算。利用§7的习题21推出：若 \(\rho\) 是 g 的非零单表示，则 \(\rho\) 具有一个重数 \(\geq2\) 的权。
\item
  对 \(\mathrm{E}_8\) 型的单代数证明同样的结果。
\end{enumerate}

\(f)\) 令 \(\mathfrak{a}\) 为可分裂单李代数。由 \(d\) )、\(e)\) 以及§7的命题7和命题8，推出下列性质等价：

\begin{enumerate}
\def\labelenumi{(\roman{enumi})}
\item
  \(\mathfrak{a}\) 具有一个非零单表示，其所有权的重数皆为 1；
\item
  \(\mathfrak{a}\) 既非 \(\mathrm{F}_4\) 型也非 \(\mathrm{E}_8\) 型。
\end{enumerate}

\exercisesection{\S~10}

\begin{enumerate}
\def\labelenumi{\arabic{enumi})}
\item
  令 \(\mathfrak{s} = \mathfrak{sl}(2, k)\)，\(\mathfrak{g} = \mathfrak{sl}(3, k)\)。把 \(\mathfrak{s}\) 与 \(\mathfrak{g}\) 的一个子代数相识别（借助 \(\mathfrak{s}\) 的 3 次不可约表示）。证明 \(\mathfrak{g}\) 的包含 \(\mathfrak{s}\) 且在 \(\operatorname{ad}_{\mathfrak{g}}\mathfrak{s}\) 下稳定的子空间都等于 \(\mathfrak{s}\) 或 \(\mathfrak{g}\)；由此推出 \(\mathfrak{s}\) 在 \(\mathfrak{g}\) 的异于 \(\mathfrak{g}\) 的子代数中是极大的。
\item
  令 \(m = \frac{1}{2}(\dim(\mathfrak{g}) + \operatorname{rk}(\mathfrak{g}))\)。g 的每个可解子代数的维数都 \(\leq m\)；若其维数为 m，则它是 Borel 子代数。（化归到代数闭的情形，并使用定理2。）
\item
  假设 \(k\) 是 \(\mathbf{R}, \mathbf{C}\)，或非离散完备超度量域。赋予格拉斯曼流形 \(\mathbf{G}(\mathfrak{g})\)（由 \(\mathfrak{g}\) 的向量子空间构成）其在 \(k\) 上的自然解析流形结构（《微分与分析流形：结果》，5.2.6）。考察 \(\mathbf{G}(\mathfrak{g})\) 中由下列对象组成的子集：
\end{enumerate}

\begin{enumerate}
\def\labelenumi{(\roman{enumi})}
\item
  子代数
\item
  可解子代数
\item
  幂零子代数
\item
  由幂零元素组成的子代数
\item
  Borel 子代数。
\end{enumerate}

证明这些子集都是闭的（对 (v)，使用习题2）。由此推出，当 \(k\) 局部紧时，这些子集都是紧的。

用例子表明，\(\mathbf{G}(\mathfrak{g})\) 中由下列对象组成的子集：

\begin{enumerate}
\def\labelenumi{(\roman{enumi})}
\setcounter{enumi}{5}
\item
  嘉当子代数
\item
  在 \(\mathfrak{g}\) 中约化的子代数
\item
  半单子代数
\item
  可分解子代数
\end{enumerate}

即使当 \(k = \mathbf{C}\) 时，也未必是闭的。

\begin{enumerate}
\def\labelenumi{\arabic{enumi})}
\setcounter{enumi}{3}
\tightlist
\item
  假设 \(k = \mathbf{C}\)。令 \(G = \operatorname{Int}(\mathfrak{g}) = \operatorname{Aut}_0(\mathfrak{g})\)，并令 \(B\) 为 \(G\) 的积分子群，其李代数 \(\mathfrak{b}\) 是 \(\mathfrak{g}\) 的 Borel 子代数。证明 \(B\) 是 \(\mathfrak{b}\) 在 \(G\) 中的正规化子（使用§3的习题4和§5的习题11）。借助习题3推出 \(G / B\) 是紧的。
\end{enumerate}

¶5) 假设 g 可分裂。若 h 是 g 的分裂嘉当子代数，则以 \(\mathrm{E}(\mathfrak{h})\) 表示 \(\mathrm{Aut}_{e}(\mathfrak{g})\) 中由诸 \(e^{\operatorname{ad} x}, x \in \mathfrak{g}^{\alpha}(\mathfrak{h}), \alpha \in \mathrm{R}(\mathfrak{g}, \mathfrak{h})\) 生成的子群，参见第VII章§3第2节。

\begin{enumerate}
\def\labelenumi{\alph{enumi})}
\item
  令 \(\mathfrak{b}\) 为 \(\mathfrak{g}\) 的包含 \(\mathfrak{h}\) 的 Borel 子代数，令 \(\mathfrak{h}_1\) 为 \(\mathfrak{b}\) 的嘉当子代数。证明 \(\mathfrak{h}_1\) 与 \(\mathfrak{h}\) 经 \(\mathrm{E}(\mathfrak{h})\) 的一个元素共轭。由此推出 \(\mathrm{E}(\mathfrak{h}) = \mathrm{E}(\mathfrak{h}_1)\)。
\item
  令 \(\mathfrak{h}'\) 为 \(\mathfrak{g}\) 的分裂嘉当子代数。证明 \(\mathrm{E}(\mathfrak{g}) = \mathrm{E}(\mathfrak{h}')\)。（若 \(\mathfrak{b}'\) 是包含 \(\mathfrak{h}'\) 的 Borel 子代数，选取 \(\mathfrak{h}_1\) 为 \(\mathfrak{b} \cap \mathfrak{b}'\) 的一个嘉当子代数，并应用 \(a\) ) 证明 \(\mathrm{E}(\mathfrak{h}) = \mathrm{E}(\mathfrak{h}_1) = \mathrm{E}(\mathfrak{h}')\)。）
\item
  令 \(x\) 为 \(\mathfrak{g}\) 的幂零元素。证明 \(e^{\mathrm{ad}x} \in \mathrm{E}(\mathfrak{h})\)。（利用 b) 和定理1的推论2，化归到 \(x \in [\mathfrak{b}, \mathfrak{b}]\) 的情形。）由此推出 \(\mathrm{E}(\mathfrak{h}) = \mathrm{Aut}_e(\mathfrak{g})\)。
\end{enumerate}

\begin{enumerate}
\def\labelenumi{\arabic{enumi})}
\setcounter{enumi}{5}
\tightlist
\item
  \begin{enumerate}
  \def\labelenumii{\alph{enumii})}
  \tightlist
  \item
    证明下列性质等价：
  \end{enumerate}
\end{enumerate}

\begin{enumerate}
\def\labelenumi{(\roman{enumi})}
\item
  g 没有 \(\neq 0\) 的幂零元素。
\item
  g 没有异于 g 的抛物子代数。
\end{enumerate}

（使用定理1的推论2。）

称这样的代数为各向异性的。

\begin{enumerate}
\def\labelenumi{\alph{enumi})}
\setcounter{enumi}{1}
\tightlist
\item
  令 \(\mathfrak{p}\) 为 \(\mathfrak{g}\) 的极小抛物子代数，\(\mathfrak{r}\) 为 \(\mathfrak{p}\) 的根基，\(\mathfrak{s} = \mathfrak{p} / \mathfrak{r}\)。证明 \(\mathfrak{s}\) 是各向异性的。（注意，若 \(\mathfrak{q}\) 是 \(\mathfrak{s}\) 的抛物子代数，则 \(\mathfrak{q}\) 在 \(\mathfrak{p}\) 中的逆像是 \(\mathfrak{g}\) 的抛物子代数，参见§3习题5 a)。）
\end{enumerate}

¶ 7) a) 证明 k 的下列性质等价：

\begin{enumerate}
\def\labelenumi{(\roman{enumi})}
\item
  每个各向异性的半单 \(k\)-李代数都化为 0。
\item
  每个半单 k-李代数都具有 Borel 子代数。
\end{enumerate}

（为证 (i) \(\Longrightarrow\) (ii)，使用习题6。）

\begin{enumerate}
\def\labelenumi{\alph{enumi})}
\setcounter{enumi}{1}
\tightlist
\item
  证明 (i) 和 (ii) 蕴含 \(^{19}\)：
\end{enumerate}

\begin{enumerate}
\def\labelenumi{(\roman{enumi})}
\setcounter{enumi}{2}
\tightlist
\item
  每个作为域的有限维 \(k\)-代数都是交换的。（或者再说一遍：\(k\) 的每个代数扩张的 Brauer 群都化为 0。）
\end{enumerate}

（使用该代数中迹为零的元素所成的李代数。）

\begin{enumerate}
\def\labelenumi{\alph{enumi})}
\setcounter{enumi}{2}
\tightlist
\item
  证明 (i) 和 (ii) 由下列性质蕴含：
\end{enumerate}

\begin{enumerate}
\def\labelenumi{(\roman{enumi})}
\setcounter{enumi}{3}
\tightlist
\item
  对任意由齐次多项式 \(f_{\alpha} \in k[(X_i)_{i \in \mathrm{I}}]\)（次数 \(\geq 1\)）组成且满足 \(\sum_{\alpha} \deg f_{\alpha} < \operatorname{Card}(\mathrm{I})\) 的有限族，存在不全为零的元素 \(x_i \in k\)，使得 \(f_{\alpha}((x_i)_{i \in \mathrm{I}}) = 0\) 对所有 \(\alpha\) 成立。
\end{enumerate}

（使用§8的命题5。）

\exercisesection{\S~11}

\begin{enumerate}
\def\labelenumi{\arabic{enumi})}
\tightlist
\item
  令 \(\mathfrak{g} = \mathfrak{sl}(2, k)\)。置 \(G = \operatorname{Aut}_e(\mathfrak{g})\)；该群可与 \(\mathbf{PSL}_2(k)\) 相识别，参见第VII章§3第1节注2。
\end{enumerate}

\begin{enumerate}
\def\labelenumi{\alph{enumi})}
\item
  \(\mathfrak{g}\) 的每个幂零元素都 G-共轭于 \(\left( \begin{array}{cc}0 & \lambda \\ 0 & 0 \end{array} \right)\)（对某个 \(\lambda \in k\)）。这样的元素是主的当且仅当它非零。
\item
  元素 \(\begin{pmatrix}0&\lambda\\0&0\end{pmatrix},\begin{pmatrix}0&\mu\\0&0\end{pmatrix}\)（其中 \(\lambda,\mu\in k^{*}\)）G-共轭当且仅当 \(\lambda^{-1}\mu\) 是 k 中的平方。
\item
  \(\mathfrak{g}\) 的每个单元素都 G-共轭于 \(\left( \begin{array}{cc}1 & 0\\ 0 & -1 \end{array} \right)\)。
\end{enumerate}

\begin{enumerate}
\def\labelenumi{\arabic{enumi})}
\setcounter{enumi}{1}
\tightlist
\item
  令 \(A = \begin{pmatrix} 0 & 1 \\ 0 & 0 \end{pmatrix}\)，\(B = \begin{pmatrix} 0 & 0 \\ 1 & 0 \end{pmatrix}\)。则 \(A\) 是幂零的，而 \(AB\) 不是，且
\end{enumerate}

\[
[ A, [ A, B ] ] = \left( \begin{array}{c c} 0 & - 2 \\ 0 & 0 \end{array} \right).
\]

由此推出，引理5不能推广到特征 2 的域上。

\begin{enumerate}
\def\labelenumi{\arabic{enumi})}
\setcounter{enumi}{2}
\tightlist
\item
  令 \(\mathfrak{r}\) 为 \(\mathfrak{g}\) 的根基，\(\mathfrak{s} = \mathfrak{g} / \mathfrak{r}\)。证明下列条件等价：
\end{enumerate}

\begin{enumerate}
\def\labelenumi{(\roman{enumi})}
\item
  g 不含 \(sl_{2}\) 三元组；
\item
  \(\mathfrak{s}\) 不含 \(\mathfrak{sl}_2\) 三元组；
\item
  \(\mathfrak{s}\) 是各向异性的（§10习题6）；
\item
  \(\mathfrak{s}\) 不含非零的可对角化元素（§3习题10）。
\end{enumerate}

（使用§3的命题2和习题10 a)。）

\begin{enumerate}
\def\labelenumi{\arabic{enumi})}
\setcounter{enumi}{3}
\tightlist
\item
  令 V 为维数 \(n \geq 2\) 的向量空间，\(\mathfrak{g} = \mathfrak{sl}(V)\)，\(G = \mathbf{PGL}(V)\) 与 g 的一个自同构群相识别。g 中的 \(sl_{2}\) 三元组赋予 V 一个忠实的 \(\mathfrak{sl}(2, k)\)-模结构，反之，任何此类结构都来自某个 \(sl_{2}\) 三元组；\(sl_{2}\) 三元组是主的当且仅当相应的 \(\mathfrak{sl}(2, k)\)-模是单模；两个 \(sl_{2}\) 三元组 G-共轭当且仅当相应的 \(\mathfrak{sl}(2, k)\)-模同构。由此推出，\(sl_{2}\) 三元组的 G-共轭类与满足下列条件的族 \((m_{1}, m_{2}, \ldots)\)（由整数 \(\geq 0\) 组成）一一对应
\end{enumerate}

\[
m _ {1} + 2 m _ {2} + 3 m _ {3} + \dots = n \quad \text { and } \quad m _ {1} <   n.
\]

¶ 5) 假设 g 是半单的。令 a 为 g 的一个子代数，它在 g 中约化、与 g 同秩，且含有 g 的一个主 \(sl_{2}\) 三元组。证明 \(a = g\)。

¶6) 假设 g 是绝对单的，并以 h 表示其 Coxeter 数。令 x 为 g 的幂零元素。证明 \((\operatorname{ad} x)^{2h-1} = 0\)，且 \((\operatorname{ad} x)^{2h-2} \neq 0\) 当且仅当 x 是主的。（化归到 g 分裂且 x 含于命题10的子代数 \(n_{+}\) 中的情形。重复命题10的证明。）

¶7) 假设 g 是半单的。令 x 为 g 的幂零元素。则 x 是主的当且仅当 x 含于 g 的唯一的 Borel 子代数中。（化归到 k 代数闭的情形。使用命题10以及§3第3节的命题10。）

\begin{enumerate}
\def\labelenumi{\arabic{enumi})}
\setcounter{enumi}{7}
\item
  半单李代数具有 \(\mathfrak{sl}_2\) 三元组当且仅当它 \(\neq 0\) 且具有 Borel 子代数。
\item
  假设 g 是半单的。令 N（分别地 P）为 g 的幂零（分别地主幂零）元素的集合。
\end{enumerate}

\begin{enumerate}
\def\labelenumi{\alph{enumi})}
\item
  证明 \(\mathrm{P}\) 是 \(\mathbf{N}\) 在 Zariski 拓扑中的开子集（使用习题6）。
\item
  假设 g 可分裂。证明 P 在 N 中稠密。（使用命题10以及§10定理1的推论2。）
\end{enumerate}

\begin{enumerate}
\def\labelenumi{\arabic{enumi})}
\setcounter{enumi}{9}
\tightlist
\item
  假设 \(\mathfrak{g}\) 是可分裂半单的。令 \((x,h,y)\) 为一个 \(\mathfrak{sl}_2\) 三元组，位于 \(\mathfrak{g}\) 中。
\end{enumerate}

\begin{enumerate}
\def\labelenumi{\alph{enumi})}
\item
  证明存在一个分裂嘉当子代数 \(\mathfrak{h}\)，它含于 \(\mathfrak{g}\) 且包含 \(h\)。（使用§3的习题10 b)。）
\item
  取 \(\mathfrak{h}\) 如 \(a\) )。证明 \(h \in \mathfrak{h}_{\mathbf{Q}}\)，且存在 \(\mathrm{R}(\mathfrak{g}, \mathfrak{h})\) 的基 B，使得 \(\alpha(h) \in \{0, 1, 2\}\) 对所有 \(\alpha \in \mathrm{B}\) 成立（参见命题5）。元素 \(x\) 属于 \(\mathfrak{g}\) 的由诸 \(\mathfrak{g}^{\alpha}\)（\(\alpha \in \mathrm{B}\)）生成的子代数。
\item
  由 a)、b) 以及 Jacobson–Morozov 定理，给出“g 的每个幂零元素都含于某个 Borel 子代数”的一个新证明（参见§10定理1的推论2）。
\end{enumerate}

¶11) 令 \((x,h,y)\) 为一个主 \(\mathfrak{sl}_2\) 三元组，位于半单李代数 \(\mathfrak{g}\) 中。赋予 \(\mathfrak{g}\) 由该三元组定义的 \(\mathfrak{sl}(2,k)\)-模结构。证明如此定义的模同构于 \(\bigoplus_{i=1}^{l} V(2k_i - 2)\)，其中 \(k_i\) 是 \(\mathfrak{g}\) 上不变多项式函数代数的特征次数。（化归到 \(\mathfrak{g}\) 是可分裂单李代数的情形。使用§8第3节定理1的推论1，以及第VI章§4习题6 c。）\(^{20}\)

¶ 12) 假设 g 是半单的。令 \(x \in g\)，并令 s（分别地 n）为 x 的半单（分别地幂零）分量。令 \(a_{x}\)（分别地 \(a_{s}\)）为 x（分别地 s）在 g 中的交换子。

\begin{enumerate}
\def\labelenumi{\alph{enumi})}
\item
  证明 \(n\) 是半单代数 \(\mathcal{D}(\mathfrak{a}_s)\) 的幂零元素，且 \(n\) 在 \(\mathfrak{a}_s\) 中的交换子等于 \(\mathfrak{a}_x\)。由此推出 \(\dim \mathfrak{a}_x < \dim \mathfrak{a}_s\)（若 \(n \neq 0\)，即若 \(x\) 不是半单的）。
\item
  证明 \(\dim \mathfrak{a}_x = \mathrm{rk}(\mathfrak{g})\) 当且仅当 \(n\) 是 \(\mathcal{D}(\mathfrak{a}_s)\) 的主幂零元素。
\item
  置 \(\mathrm{G} = \mathrm{Aut}_{e}(\mathfrak{g})\)。证明对所有 \(\lambda \in k\)，存在 \(\sigma_{\lambda} \in G\) 使得 \(\sigma_{\lambda}x = s + \lambda^{2}n\)。（证明：若 \(n \neq 0\)，则存在一个 \(sl_{2}\) 三元组（在 \(a_{s}\) 中），其第一分量为 n，由此得到一个同态 \(\varphi : \mathbf{SL}(2, k) \to \mathbf{G}\)；取 \(\sigma_{\lambda}\) 为在 \(\varphi\) 下 \(\mathbf{SL}(2, k)\) 的适当对角元素的像。）由此推出 s 属于 G.x 在 Zariski 拓扑中的闭包。
\end{enumerate}

\(d\) ) 证明，若 \(x\) 不是半单的，则 \(x\) 不属于 G.s 在 Zariski 拓扑中的闭包（利用不等式 \(\dim \mathfrak{a}_x < \dim \mathfrak{a}_s\)，参见 \(a\) )）。

\begin{enumerate}
\def\labelenumi{\alph{enumi})}
\setcounter{enumi}{4}
\tightlist
\item
  假设 \(k\) 代数闭。证明下列性质等价：
\end{enumerate}

\begin{enumerate}
\def\labelenumi{(\roman{enumi})}
\item
  \(x\) 是半单的；
\item
  G.x 在 g 中是 Zariski 拓扑意义下闭的。
\end{enumerate}

（蕴含关系 (ii) \(\Longrightarrow\) (i) 由 c) 得到。若 (i) 成立且 \(x'\) 属于 G.x 的闭包，则§8的习题15表明 \(s'\)（\(x'\) 的半单分量）属于 G.x，于是 \(x'\) 属于 G. \(s'\) 的闭包；对 \(x'\) 与 \(s'\) 应用 d) 即得结论。）

\(f\) ) 假设 \(k\) 代数闭。令 \(\mathrm{F}_x\) 为由元素 \(y\in \mathfrak{g}\)（满足 \(f(x) = f(y)\)，对每个不变多项式函数 \(f\)（\(\mathfrak{g}\) 上的））组成的集合；我们有 \(y\in \mathrm{F}_x\) 当且仅当 \(y\) 的半单分量 G-共轭于 \(s\)（§8习题17）。证明 \(\mathrm{F}_x\) 是 G 的有限多条轨道的并，且轨道条数 \(\leq 3^{l(x)}\)，其中 \(l(x)\) 是 \(\mathcal{D}(\mathfrak{a}_s)\) 的秩。这些轨道中只有一条是闭的：\(s\) 的轨道；只有一条在 \(\mathrm{F}_x\) 中是开的：由元素 \(y\in \mathrm{F}_x\)（满足 \(\dim \mathfrak{a}_y = \operatorname {rk}(\mathfrak{g})\)）组成的那条。

\begin{enumerate}
\def\labelenumi{\arabic{enumi})}
\setcounter{enumi}{12}
\tightlist
\item
  假设 \(\mathfrak{g}\) 是半单的。
\end{enumerate}

\begin{enumerate}
\def\labelenumi{\alph{enumi})}
\item
  令 \((x, h, y)\) 为 g 中的主 \(sl_{2}\) 三元组，并令 b 为包含 x 的 Borel 子代数（习题7）。证明 b 含于 Im ad x。
\item
  假设 k 代数闭。证明对 g 中任意元素 z，存在 \(x, t \in g\)，其中 x 为主幂零元，使得 \(z = [x, t]\)（对包含 z 的一个 Borel 子代数应用 a)）。
\end{enumerate}

\begin{enumerate}
\def\labelenumi{\arabic{enumi})}
\setcounter{enumi}{13}
\tightlist
\item
  假设 \(\mathfrak{g}\) 是半单的。令 \(\mathfrak{p}\) 为 \(\mathfrak{g}\) 的抛物子代数，并令 \(f_{1}\)、\(f_{2}\) 为从 \(\mathfrak{g}\) 到某个有限维李代数的两个同态。证明 \(f_{1}|\mathfrak{p} = f_{2}|\mathfrak{p}\) 蕴含 \(f_{1} = f_{2}\)。（化归到 \(\mathfrak{g}\) 分裂的情形，再化归到 \(\mathfrak{g} = \mathfrak{sl}(2,k)\) 的情形，并使用第1节的引理1。）
\end{enumerate}

¶ 15) 假设 g 是半单的。

\begin{enumerate}
\def\labelenumi{\alph{enumi})}
\item
  令 \(x\) 为 \(\mathfrak{g}\) 的幂零元素。证明 \(x\) 含于 \(\operatorname{Im}(\operatorname{ad} x)^2\)（使用命题2）。由此推出 \((\operatorname{ad} x)^2 = 0\) 蕴含 \(x = 0\)。
\item
  令 \((x,h,y)\) 为一个 \(\mathfrak{sl}_2\) 三元组（在 \(\mathfrak{g}\) 中），并令 \(\mathfrak{s} = kx\oplus kh\oplus ky\)。证明下列条件等价：
\end{enumerate}

\begin{enumerate}
\def\labelenumi{(\roman{enumi})}
\item
  \(\operatorname{Im}(\operatorname{ad} x)^{2}=k.x;\)
\item
  \(\mathfrak{s}\)-模 \(\mathfrak{g} / \mathfrak{s}\) 是维数为 1 或 2 的单模之和；
\item
  \(\mathrm{ad}_{\mathfrak{g}}h\) 的异于 0、1 和 -1 的特征值只有 2 和 -2，且它们的重数都等于 1。
\end{enumerate}

\begin{enumerate}
\def\labelenumi{\alph{enumi})}
\setcounter{enumi}{2}
\item
  假设 g 是可分裂单李代数。令 h 为 g 的分裂嘉当子代数，B 为 \(\mathrm{R}(\mathfrak{g},\mathfrak{h})\) 的基，\(\gamma\) 为 \(\mathrm{R}(\mathfrak{g},\mathfrak{h})\) 相对于 B 的最高根。令 \((x,h,y)\) 为一个 \(sl_{2}\) 三元组，满足 \(h\in h_{Q}\) 且 \(\alpha(h)\geq0\) 对所有 \(\alpha\in B\)（参见命题5）。证明 b) 中的条件 (i)、(ii)、(iii) 成立当且仅当 \(h=H_{\gamma}\)，此时 \(x\in g^{\gamma}\) 且 \(y\in g^{-\gamma}\)。
\item
  保留 c) 的假设，并置 \(\mathrm{G} = \mathrm{Aut}_{0}(\mathfrak{g})\)。证明满足 (i)、(ii) 和 (iii) 的 \(sl_{2}\) 三元组是 G-共轭的（使用习题10）。证明：若 \(x\) 是 \(\mathfrak{g}\) 的满足 (i) 的非零幂零元素，则 \(G.x\) 在 Zariski 拓扑中的闭包等于 \(\{0\} \cup G.x\)。
\end{enumerate}

\begin{enumerate}
\def\labelenumi{\arabic{enumi})}
\setcounter{enumi}{15}
\tightlist
\item
  令 \((x,h,y)\) 为 g 中的一个 \(sl_{2}\) 三元组。证明：若 -2 是 k 中的平方，则元素 x-y 与 h 经 \(\operatorname{Aut}_{e}(\mathfrak{g})\) 的一个元素共轭（化归到 \(\mathfrak{g}=\mathfrak{sl}(2,k)\) 的情形）。由此推出：若 g 半单，则 x-y 是半单的，且它是正则的当且仅当 h 是正则的。
\end{enumerate}

¶ 17) 令 \((x, h, y)\) 为一个主 \(\mathfrak{sl}_2\) 三元组，位于半单代数 \(\mathfrak{g}\) 中。对所有 \(i \in \mathbf{Z}\)，以 \(\mathfrak{g}_i\) 表示 ad \(h\) 相应于特征值 \(i\) 的特征子空间；我们有 \(\mathfrak{g} = \bigoplus_{i \in \mathbf{Z}} \mathfrak{g}_i\)，且 \(\mathfrak{g}_i = 0\)（当 \(i\) 为奇数）。直和 \(\mathfrak{b}\)——诸 \(\mathfrak{g}_i\)（\(i \geq 0\)）之和——是 \(\mathfrak{g}\) 的 Borel 子代数，而 \(\mathfrak{g}_0\) 是嘉当子代数。

\begin{enumerate}
\def\labelenumi{\alph{enumi})}
\item
  证明对所有 \(z \in \mathfrak{b}\)，\(y + z\) 在 \(\mathfrak{g}\) 中的交换子的维数为 \(l = \operatorname{rk}(\mathfrak{g})\)。
\item
  令 I 为 \(\mathfrak{g}\) 上的不变多项式函数代数，\(\mathrm{P}_1, \ldots, \mathrm{P}_l\) 为生成 I 的 I 中齐次元素；置 \(\deg(\mathrm{P}_i) = k_i = m_i + 1\)。证明 \(\mathfrak{c}\)（\(x\) 在 \(\mathfrak{g}\) 中的交换子）具有基 \(x_1, \ldots, x_l\)，其中 \(x_i \in \mathfrak{g}_{2m_i}\)（参见习题11）。
\item
  令 \(i \in \{1, l\}\)，并令 \(J_i\)（分别地 \(K_i\)）为由 \(j \in (1, l)\)（满足 \(m_j = m_i\)（分别地 \(m_j < m_i\)））组成的集合。令 \(f_i \in k[X_1, \ldots, X_l]\) 为使得
\end{enumerate}

\[
f _ {i} \left(a _ {1}, \dots , a _ {l}\right) = \mathrm{P} _ {i} \left(y + \sum_ {j = 1} ^ {j = l} a _ {j} x _ {j}\right) \quad \text { for } \left(a _ {j}\right) \in k ^ {l}.
\]

证明 \(f_{i}\) 是线性型 \(\mathrm{L}_i\)（关于诸 \(\mathrm{X}_j\)，\(j \in \mathrm{J}_i\)）与一个多项式（关于诸 \(\mathrm{X}_j\)，\(j \in \mathrm{K}_i\)）之和。（若 \(t \in k^*\)，则 \(\mathfrak{g}\) 上等于 \(t^i\)（在 \(\mathfrak{g}_i\) 上）的自同构属于 \(\operatorname{Aut}_0(\mathfrak{g})\)，它把 \(y\) 变为 \(t^{-2}y\)，把 \(x_j\) 变为 \(t^{2m_j}x_j\)。利用 \(\mathrm{P}_i\) 在该自同构下的不变性。）

\(d)\) 令 \(\mathrm{P}\) 为从 \(\mathfrak{g}\) 到 \(k^l\) 的、由诸 \(\mathrm{P}_i\) 定义的映射。若 \(z \in \mathfrak{g}\)，以 \(\mathrm{D}_z\mathrm{P} : \mathfrak{g} \to k^l\) 表示 \(\mathrm{P}\) 在 \(z\) 处的切线性映射（第VII章附录I第2节）。

证明 \(\mathfrak{c}\cap\operatorname{Im}\operatorname{ad}(y-x)=0\)（相对于由所给 \(sl_{2}\) 三元组生成的子代数，把 g 分解为单子模的直和）。由此推出 \(D_{y-x}P\) 在 c 上的限制是从 c 到 \(k^{l}\) 的同构（利用前一习题及§8的习题18 a)）。利用这一结果证明，c) 中定义的线性型 \(L_{i}\) 的行列式 \(\neq0\)，并由此推出下列结果：

\(d_{1})\) 多项式 \(f_{1},\ldots ,f_{l}\) 代数无关，且生成 \(k[\mathrm{X}_1,\dots ,\mathrm{X}_l]\)；

\(d_{2})\) 映射 \(z \mapsto \mathrm{P}(y + z)\)（从 \(\mathfrak{c}\) 到 \(k^{l}\)）是双射多项式映射，且其逆映射也是多项式映射；

\(d_{3})\) 对所有 \(z\in y + \mathfrak{c}\)，线性映射 \(\mathrm{D}_z\mathrm{P}|\mathfrak{c}\) 的秩为 \(l\)。

特别地，映射 \(\mathrm{P}:\mathfrak{g}\to k^l\) 是满射。

\begin{enumerate}
\def\labelenumi{\alph{enumi})}
\setcounter{enumi}{4}
\tightlist
\item
  若 \(k\) 代数闭，则 \(\mathfrak{g}\) 中每个交换子维数为 \(l\) 的元素都经 \(\mathrm{Aut}_0(\mathfrak{g})\) 共轭于 \(y + \mathfrak{c}\) 的唯一元素。
\end{enumerate}

\(f\) ) 给出一个没有主 \(\mathfrak{sl}_2\) 三元组且映射 \(\mathrm{P}\) 非满射的单李代数的例子（取 \(k = \mathbf{R}\) 及 \(l = 1\)）。

\exercisesection{\S~13}

\begin{enumerate}
\def\labelenumi{\arabic{enumi})}
\tightlist
\item
  可分裂单李代数的 \(\leq 80\) 的维数为：
\end{enumerate}

3 (A \(_{1}\) = B \(_{1}\) = C \(_{1}\) ), 8 (A \(_{2}\) ), 10 (B \(_{2}\) = C \(_{2}\) ), 14 (G \(_{2}\) ), 15 (A \(_{3}\) = D \(_{3}\) ), 21 (B \(_{3}\) 与 C \(_{3}\) ), 24 (A \(_{4}\) ), 28 (D \(_{4}\) ), 35 (A \(_{5}\) ), 36 (B \(_{4}\) 与 C \(_{4}\) ), 45 (D \(_{5}\) ), 48 (A \(_{6}\) ), 52 (F \(_{4}\) ), 55 (B \(_{5}\) 与 C \(_{5}\) ), 63 (A \(_{7}\) ), 66 (D \(_{6}\) ), 78 (B \(_{6}\) , C \(_{6}\) 与 E \(_{6}\) ), 80 (A \(_{8}\) ).

\begin{enumerate}
\def\labelenumi{\arabic{enumi})}
\setcounter{enumi}{1}
\tightlist
\item
  令 \((\mathfrak{g},\mathfrak{h})\) 为分裂单李代数，B 为 \(R(\mathfrak{g},\mathfrak{h})\) 的基。维数最小的单 \(\mathfrak{g}\)-模 \(E(\lambda)\)（\(\lambda \in P_{++} - \{0\}\)）是那些 \(\lambda\) 为下列权之一的模：
\end{enumerate}

\(\varpi_{1}\)（A \(_{1}\)）；\(\varpi_{1}\) 与 \(\varpi_{l}\)（A \(_{l}\)，\(l \geq 2\)）；\(\varpi_{1}\)（B \(_{l}\) 与 C \(_{l}\)，\(l \geq 2\)）；\(\varpi_{1}\)、\(\varpi_{3}\) 与 \(\varpi_{4}\)（D \(_{4}\)）；\(\varpi_{1}\)（D \(_{l}\)，\(l \geq 5\)）；\(\varpi_{1}\) 与 \(\varpi_{6}\)（E \(_{6}\)）；\(\varpi_{7}\)（E \(_{7}\)）；\(\varpi_{8}\)（E \(_{8}\)）；\(\varpi_{4}\)（F \(_{4}\)）；\(\varpi_{1}\)（G \(_{2}\)）。任何两个这样的模都可经 g 的一个自同构相互变换。

\(E_{8}\) 型是唯一使伴随表示达到最小维数的类型。

\begin{enumerate}
\def\labelenumi{\arabic{enumi})}
\setcounter{enumi}{2}
\tightlist
\item
  \begin{enumerate}
  \def\labelenumii{\alph{enumii})}
  \tightlist
  \item
    定义一个从 \(\mathfrak{sl}(4,k)\) 到正交李代数 \(\mathfrak{o}_S(6,k)\) 的同构。（利用第1节.V 的表示 \(\bigwedge^2\sigma\) 是正交的且维数为 6 这一事实。）\(\mathfrak{sl}(4,k)\) 的两个次数为 4 的不可约表示类型对应于 \(\mathfrak{o}_S(6,k)\) 的两个半旋量表示。
  \end{enumerate}
\end{enumerate}

\begin{enumerate}
\def\labelenumi{\alph{enumi})}
\setcounter{enumi}{1}
\item
  定义一个从 \(\mathfrak{sp}(4,k)\) 到正交李代数 \(\mathfrak{o}_S(5,k)\) 的同构。（利用第3节.V 的表示 \(\sigma_2\) 是正交的且维数为 5 这一事实。）\(\mathfrak{sp}(4,k)\) 的次数为 4 的不可约表示对应于 \(\mathfrak{o}_S(5,k)\) 的旋量表示。
\item
  定义一个同构 \(\mathfrak{sl}(2,k)\times\mathfrak{sl}(2,k)\to\mathfrak{o}_{S}(4,k)\)（利用两个 \(\mathfrak{sl}(2,k)\) 因子的恒等表示的张量积）。借助第I章§6习题26重新得到这一结果。
\end{enumerate}

\begin{enumerate}
\def\labelenumi{\arabic{enumi})}
\setcounter{enumi}{3}
\tightlist
\item
  设 \(S\) 为 \(n\) 阶方阵
\end{enumerate}

\[
(\delta_ {i, n + 1 - i}) = \left( \begin{array}{c c c c c} 0 & 0 & \dots & 0 & 1 \\ 0 & 0 & \dots & 1 & 0 \\ \vdots & \vdots & \ddots & \vdots & \vdots \\ 0 & 1 & \dots & 0 & 0 \\ 1 & 0 & \dots & 0 & 0 \end{array} \right).
\]

\(\mathfrak{o}_{S}(n,k)\) 的元素是关于副对角线反对称的矩阵 \((a_{ij})\)：

\[
a _ {i, j} = - a _ {n + 1 - j, n + 1 - i} \quad \text {   for   every   pair   } (i, j).
\]

李代数 \(\mathfrak{o}_S(n,k)\) 是 \(\mathrm{D}_{n/2}\) 型的可分裂单李代数（若 \(n\) 为偶数且 \(\geq 6\)），而是 \(\mathrm{B}_{(n-1)/2}\) 型（若 \(n\) 为奇数且 \(\geq 5\)）。\(\mathfrak{o}_S(n,k)\) 的对角（分别地，上三角）元素构成分裂嘉当（分别地，Borel）子代数。

\begin{enumerate}
\def\labelenumi{\arabic{enumi})}
\setcounter{enumi}{4}
\item
  记号沿用第1节.IV，\(\mathrm{A}_l\) 型。证明：若 \(n\) \(\geq 0\)，则 \(\mathbf{S}^n\sigma\) 是 \(\mathfrak{sl}(l + 1,k)\) 的最高权为 \(n\varpi_1\) 的不可约表示。（把 \(\mathbf{S}^n\sigma\) 实现在次数为 \(n\) 的齐次多项式（关于 \(\mathrm{X}_0,\dots ,\mathrm{X}_l\)）的空间上，并注意在相差一个位似的意义下，唯一的多项式 \(f\)（满足 \(\partial f / \partial \mathrm{X}_i = 0\)，对 \(i\geq 1\)）是 \(\mathrm{X}_0^n\)。）证明这个表示的所有权都是重数 1 的。
\item
  记号沿用第2节.IV，\(\mathrm{B}_l\) 型。证明：若 \(1 \leq r \leq l - 1\)，则 \(\operatorname{E}(\varpi_r)\) 的维数是 \(\binom{2l + 1}{r}\)，并由此给出 \(\bigwedge^r \sigma\) 是最高权为 \(\varpi_r\) 的基本表示这一事实的另一个证明。
\item
  记号沿用第2节.VII，\(\mathrm{B}_l\) 型。证明 \(\mathbf{O}_0^+ (\Psi)\) 是 \(\mathbf{SO}(\Psi)\) 的换位子群。（注意 \(\mathbf{O}(\Psi)\) 等于 \(\{\pm 1\} \times \mathbf{SO}(\Psi)\)，因而与 \(\mathbf{SO}(\Psi)\) 有相同的换位子群，然后应用《代数学》第IX章§9习题11 b）。）由此推出 \(\operatorname{Aut}_e(\mathfrak{g}) = \mathbf{O}_0^+ (\Psi)\)。
\item
  记号沿用第3节.IV，\(C_{l}\) 型（\(l \geq 1\)）。证明 \(S^{2}\sigma\) 等价于 g 的伴随表示。
\item
  记号沿用第3节.VII，\(C_l\) 型（\(l \geq 1\)）。特别地，我们把 \(\mathrm{Aut}_e(\mathfrak{g})\) 与 \(\mathbf{Sp}(\varPsi)/\{\pm 1\}\) 的一个子群等同起来。证明辛平延（《代数学》第IX章§4习题6）在 \(\mathbf{Sp}(\varPsi)/\{\pm 1\}\) 中的像属于 \(\mathrm{Aut}_e(\mathfrak{g})\)。由此推出 \(\mathrm{Aut}_e(\mathfrak{g}) = \mathbf{Sp}(\varPsi)/\{\pm 1\}\)（《代数学》第IX章§5习题11），并且 \(\mathrm{Aut}(\mathfrak{g})/\mathrm{Aut}_e(\mathfrak{g})\) 可以与 \(k^*/k^{*2}\) 等同起来。
\end{enumerate}

¶ 10) 记号沿用第4节.IV，D \(_{l}\) 型（l ≥ 2）。

\begin{enumerate}
\def\labelenumi{\alph{enumi})}
\tightlist
\item
  令 \(x\) 与 \(y\) 为 \(\bigwedge^{l} \mathrm{V}\) 中由下式定义的元素
\end{enumerate}

\[
x = e _ {1} \wedge \dots \wedge e _ {l - 1} \wedge e _ {l} \quad \text { and } \quad y = e _ {1} \wedge \dots \wedge e _ {l - 1} \wedge e _ {- l}.
\]

元素 x 是权为 \(2\varpi_{l}\) 的本原元，y 是权为 \(2\varpi_{l-1}\) 的本原元。\(\bigwedge^{l}V\) 中由 x（分别地，y）生成的子模 X（分别地，Y）同构于 \(\mathrm{E}(2\varpi_{l})\)（分别地，\(\mathrm{E}(2\varpi_{l-1})\)）。通过计算维数证明 \(\bigwedge^{l}V = X \oplus Y\)；特别地，\(\bigwedge^{l}V\) 是两个互不同构的单模之和。

\begin{enumerate}
\def\labelenumi{\alph{enumi})}
\setcounter{enumi}{1}
\tightlist
\item
  令 \(e = e_{1} \wedge \cdots \wedge e_{l} \wedge e_{-1} \wedge \cdots \wedge e_{-l} \in \bigwedge^{2l}V\)，并令 \(\varPsi_{l}\) 为 \(\varPsi\) 到 \(\bigwedge^{l}V\) 上的扩张。设 \(z \in \bigwedge^{l}V\)。证明下列等价关系：
\end{enumerate}

\[
z \in X \Longleftrightarrow z \wedge t = \Psi_ {l} (z, t) e \quad \text {   for   all   } t \in \bigwedge^ {l} V
\]

\[
z \in \mathrm{Y} \Longleftrightarrow z \wedge t = - \Psi_ {l} (z, t) e \quad \text {   for   all   } t \in \bigwedge^ {l} \mathrm{V}.
\]

（若以 \(X'\) 与 \(Y'\) 表示由右边定义的子空间，先证明 \(X'\) 与 \(Y'\) 在 \(\mathfrak{g}\) 作用下稳定，且分别含有 \(x\) 与 \(y\)。）

\begin{enumerate}
\def\labelenumi{\alph{enumi})}
\setcounter{enumi}{2}
\item
  假设 z 是纯元素（《代数学》第III章§11第13节），并以 \(M_{z}\) 表示与之相伴的 V 的 l 维子空间。证明 \(M_{z}\) 全迷向当且仅当 z 属于 X 或 Y。（利用下述事实：当 \(M_{z}\) 全迷向时，存在一个正交变换把 \(M_{x}\) 变到它；当 \(M_{z}\) 不全迷向时，构造一个 l-向量 t 使得 \(z \wedge t = 0, \Psi_{l}(z, t) = 1\)，并应用上面的 b）。）当 \(z \in X\)（分别地，\(z \in Y\)）时，\(\mathrm{M}_{x}/(\mathrm{M}_{x} \cap \mathrm{M}_{z})\) 的维数是偶（分别地，奇）整数，参见《代数学》第IX章§6习题18 d）。
\item
  设 s 是 V 的一个正向（分别地，反向）相似变换。证明 \(\bigwedge^{l}s\) 保持 X 与 Y 稳定（分别地，互换 X 与 Y）。
\end{enumerate}

\begin{enumerate}
\def\labelenumi{\arabic{enumi})}
\setcounter{enumi}{10}
\item
  记号沿用第4节.VII，\(\mathrm{D}_l\) 型（\(l \geq 3\)）。证明 \(\mathbf{O}_0^+(\varPsi)\) 是 \(\mathbf{SO}(\varPsi)\) 的换位子群。（应用《代数学》第IX章§6习题17 b）及《代数学》第IX章§9习题11 b）。）由此推出 \(\mathrm{Aut}_e(\mathfrak{g})\) 等于 \(\mathbf{O}_0^+(\varPsi)\) 在 \(\mathbf{SO}(\varPsi)/\{\pm 1\}\) 中的像。此外，\(-1\) 属于 \(\mathbf{O}_0^+(\varPsi)\) 当且仅当 \(l\) 为偶数或 \(-1\) 是 \(k\) 中的平方（《代数学》第IX章§9习题11 c））。
\item
  \(\mathfrak{sl}(n,k)\) 的 Killing 型是 \((X,Y)\mapsto 2n\operatorname {Tr}(XY)\)。\(\mathfrak{sp}(n,k)\)（\(n\) 为偶数时）的 Killing 型是 \((X,Y)\mapsto (n + 2)\operatorname {Tr}(XY)\)。\(\mathfrak{o}_S(n,k)\)（其中 \(S\) 是秩为 \(n\) 的非退化对称矩阵）的 Killing 型是 \((X,Y)\mapsto (n - 2)\operatorname {Tr}(XY)\)。
\item
  \(\mathfrak{g}\) 的不变多项式函数代数由下列元素生成：
\end{enumerate}

\(a\) ) 在 \(\mathrm{A}_l\) 型情形，由函数 \(X\mapsto \mathrm{Tr}(X^i)\)（\(2\leq i\leq l + 1\)）生成；

\begin{enumerate}
\def\labelenumi{\alph{enumi})}
\setcounter{enumi}{1}
\item
  在 \(B_{l}\) 型情形，由函数 \(X \mapsto \operatorname{Tr}(X^{2i})\)（\(1 \leq i \leq l\)）生成；
\item
  在 \(C_{l}\) 型情形，由函数 \(X \mapsto \operatorname{Tr}(X^{2i})\)（\(1 \leq i \leq l\)）生成；
\end{enumerate}

\(d\) ) 在 \(\mathrm{D}_l\) 型情形，由函数 \(X \mapsto \operatorname{Tr}(X^{2i})\)（\(1 \leq i \leq l - 1\)）以及两个多项式函数 \(\tilde{f}\) 之一（满足 \(\tilde{f}(X)^2 = (-1)^l \det(X)\)）生成。

\begin{enumerate}
\def\labelenumi{\arabic{enumi})}
\setcounter{enumi}{13}
\tightlist
\item
  \begin{enumerate}
  \def\labelenumii{\alph{enumii})}
  \tightlist
  \item
    令 G 为按 §7 习题26 的程序与 \(\mathfrak{g} = \mathfrak{sl}(n, k)\) 相伴的群。\(k^{n}\) 上的自然 g-模结构给出一个同态 \(\varphi : G \to \mathbf{GL}(n, k)\)。利用该处 h) 证明 \(\varphi\) 是单射，并利用该处 f) 证明
  \end{enumerate}
\end{enumerate}

\[
\operatorname{Im} (\varphi) = \mathbf {S L} (n, k).
\]

\begin{enumerate}
\def\labelenumi{\alph{enumi})}
\setcounter{enumi}{1}
\item
  令 E 为有限维 \(\mathfrak{sl}(n,k)\)-模，\(\rho\) 为 \(\mathfrak{sl}(n,k)\) 的相应表示。证明存在唯一的表示 \(\pi:\mathbf{SL}(n,k)\to\mathbf{GL}(\mathrm{E})\)，使得 \(\pi(e^{x})=e^{\rho(x)}\) 对 \(\mathfrak{sl}(n,k)\) 的每个幂零元 x 成立（利用 a)）。我们称 \(\rho\) 与 \(\pi\) 相容。把 §1 第4节中对 n=2 证明的结果加以推广。
\item
  假设 \(k\) 是 \(\mathbf{R}, \mathbf{C}\)，或是一个关于离散赋值完备且剩余域特征 \(\neq 0\) 的域。证明 \(\rho\) 与 \(\pi\) 相容当且仅当 \(\pi\) 是李群同态且 \(L(\pi) = \rho\)（方法同 §1 习题18 b)）。
\end{enumerate}

\(d\) ) 对 \(\mathfrak{sp}(2n,k)\) 与 \(\mathbf{Sp}(2n,k)\) 证明类似的结果。

¶15) 令 V 为有限维 \(\geq 2\) 的向量空间，g 为 End(V) 的李子代数，\(\theta\) 为 g 的元素。我们作如下假设：

\begin{enumerate}
\def\labelenumi{(\roman{enumi})}
\item
  V 是半单 g-模；
\item
  \(\theta\) 的秩为 1（即 \(\dim \operatorname{Im}(\theta) = 1\)）；
\item
  直线 \(\operatorname{Im}(\theta)\) 生成 U(g)-模 V。
\end{enumerate}

\begin{enumerate}
\def\labelenumi{\alph{enumi})}
\item
  证明这些假设成立（对适当选取的 \(\theta\)）：当 \(\mathfrak{g} = \mathfrak{sl}(V)\)、当 \(\mathfrak{g} = \mathfrak{gl}(V)\)，或当 V 上存在非退化交错双线性形式 \(\varPsi\) 使得 \(\mathfrak{g} = \mathfrak{sp}(\varPsi)\) 或 \(\mathfrak{g} = k.1 \oplus \mathfrak{sp}(\varPsi)\) 时。在每种情形，都可取 \(\theta\) 为幂零元；仅在第二种情形（且仅在该情形）可取 \(\theta\) 为半单元。
\item
  我们将证明上述四种情形是仅有的可能情形。立即可化为 k 代数闭的情形。证明此时 V 是单 g-模，且 \(g = c \oplus s\)，其中 s 半单，c = 0 或 c = k.1；证明 V 不同构于维数 \(\geq 2\) 的 g-模的张量积；由此推出 s 是单的。
\item
  取定 \(\mathfrak{h}\) 为 \(\mathfrak{s}\) 的一个嘉当子代数，以及 \(\mathrm{R}(\mathfrak{s},\mathfrak{h})\) 的基 B。令 \(\lambda\) 为 \(\mathfrak{s}\)-模 V 的（关于 B 的）最高权，并令 \(e\) 为 V 中权为 \(\lambda\) 的非零元素。对偶模 \(\mathrm{V}^*\) 的最高权是 \(\lambda^{*} = -w_{0}\lambda\)（§7 第5节）；令 \(e^*\) 为 \(\mathrm{V}^*\) 中权为 \(\lambda^*\) 的非零元素。按惯例把 \(\mathrm{V} \otimes \mathrm{V}^*\) 与 \(\operatorname{End}(\mathrm{V})\) 等同起来。证明存在 \(x \in \mathrm{V}, y \in \mathrm{V}^*\)，使得 \(x \otimes y \in \mathfrak{g}\) 且 \(\langle x, e^* \rangle \neq 0, \langle e, y \rangle \neq 0\)（取 \(\theta\) 在 \(e^n\) 下的共轭，其中 \(n\) 是 \(\mathfrak{g}\) 的一个适当的幂零元）。利用 \(\mathfrak{g}\) 是 \(\mathfrak{h}\)-子模（\(\mathrm{V} \otimes \mathrm{V}^*\) 中的）这一事实，推出 \(\mathfrak{g}\) 含有 \(e \otimes e^*\)。由此推出 \(\lambda + \lambda^{*} = \tilde{\alpha}\)，其中 \(\tilde{\alpha}\) 是 \(\mathfrak{s}\) 的最高根。
\end{enumerate}

\(d\) ) 证明 \(\mathfrak{s}\) 不是 \(\mathrm{B}_l\) 型（\(l\geq 2\)）、\(\mathrm{D}_l\) 型（\(l\geq 4\)）、\(\mathrm{E}_6,\mathrm{E}_7,\mathrm{E}_8,\mathrm{F}_4,\mathrm{G}_2\) 型（因为由第VI章的表，\(\tilde{\alpha}\) 将是基本权，从而不可能具有上面 \(\lambda +\lambda^{*}\) 的形式）。由此推出 \(\mathfrak{s}\) 是 \(\mathrm{A}_l\) 型或 \(\mathrm{C}_l\) 型，并且在第一种情形 \(\lambda = \varpi_1\) 或 \(\varpi_1 = \varpi_1^*\)，在第二种情形 \(\tilde{\alpha} = 2\varpi_1 = \varpi_1 + \varpi_1^*\)；由于 \(\mathfrak{c} = 0\) 或 \(k.1\)，这确实给出 \(a)^{21}\) 中的四种可能。

¶16) 令 g 为 \(A_{l}\) 型（\(l \geq 2\)）的绝对单李代数，\(\bar{k}\) 为 k 的代数闭包，\(\pi : \operatorname{Gal}(\bar{k}/k) \to \operatorname{Aut}(\bar{\mathbf{R}},\bar{\mathbf{B}})\) 为 §5 习题8 中定义的同态。

\begin{enumerate}
\def\labelenumi{\alph{enumi})}
\tightlist
\item
  假设 \(\pi\) 平凡。证明存在恰好两个双边理想 \(\mathfrak{m}\) 与 \(\mathfrak{m}'\)（\(\mathrm{U}(\mathfrak{g})\) 的），使得 \(\mathrm{D} = \mathrm{U}(\mathfrak{g}) / \mathfrak{m}\) 与 \(\mathrm{D}' = \mathrm{U}(\mathfrak{g}) / \mathfrak{m}'\) 都是维数为 \((l + 1)^2\) 的中心单代数（利用 §7 习题8）。\(\mathrm{U}(\mathfrak{g})\) 的主反自同构互换 \(\mathfrak{m}\) 与 \(\mathfrak{m}'\)；特别地，\(\mathrm{D}'\) 同构于 D 的反代数。复合映射 \(\mathfrak{g} \to \mathrm{U}(\mathfrak{g}) \to \mathrm{D}\) 把 \(\mathfrak{g}\) 与 D 中由迹为零的元素组成的李子代数 \(\mathfrak{sl}_{\mathrm{D}}\) 等同起来。此外，\(\mathfrak{g}\) 可分裂（从而同构于 \(\mathfrak{sl}(l + 1,k)\)）当且仅当 D 同构于 \(\mathbf{M}_{l + 1}(k)\)。
\end{enumerate}

反之，若 \(\Delta\) 是维数为 \((l + 1)^2\) 的中心单代数，则李代数 \(\mathfrak{sl}_{\Delta}\) 是 \(A_l\) 型的绝对单李代数，且相应的同态 \(\pi\) 平凡。两个这样的代数 \(\mathfrak{sl}_{\Delta}\) 与 \(\mathfrak{sl}_{\Delta'}\) 同构当且仅当 \(\Delta\) 与 \(\Delta'\) 同构或反同构。

\begin{enumerate}
\def\labelenumi{\alph{enumi})}
\setcounter{enumi}{1}
\tightlist
\item
  假设 \(\pi\) 非平凡。由于 \(\operatorname{Aut}(\bar{\mathbf{R}},\bar{\mathbf{B}})\) 有两个元素，\(\pi\) 的核是 \(\operatorname{Gal}(\bar{k}/k)\) 的指标为 2 的开子群，按 Galois 理论它对应于 k 的一个二次扩张 \(k_{1}\)。把 \(k_{1}\) 的非平凡对合记作 \(x\mapsto\bar{x}\)。
\end{enumerate}

证明存在唯一的双边理想 \(\mathfrak{m}\)（\(\mathrm{U}(\mathfrak{g})\) 的），使得 \(\mathrm{D} = \mathrm{U}(\mathfrak{g}) / \mathfrak{m}\) 是维数为 \(2(l + 1)^2\) 的单代数，且其中心是 \(k\) 的二次扩张（方法同上）。D 的中心可与 \(k_{1}\) 等同起来。理想 \(\mathfrak{m}\) 在 \(\mathrm{U}(\mathfrak{g})\) 的主反自同构下稳定；于是过渡到商即定义出 D 的一个对合反自同构 \(\sigma\)，使得 \(\sigma(x) = \bar{x}\) 对所有 \(x \in k_{1}\) 成立。复合映射 \(\mathfrak{g} \to \mathrm{U}(\mathfrak{g}) \to \mathrm{D}\) 把 \(\mathfrak{g}\) 与李子代数 \(\mathfrak{s}\mathfrak{u}_{\mathrm{D},\sigma}\)——由元素 \(x\)（满足 \(\sigma(x) = -x\) 且 \(\mathrm{Tr}_{\mathrm{D}/k_1}(x) = 0\)）组成——等同起来。反之，若 \(\Delta\) 是中心单 \(k_{1}\)-代数（维数为 \((l + 1)^2\)），且带有一个对合反自同构 \(\sigma\)（它在 \(k_{1}\) 上的限制为 \(x \mapsto \bar{x}\)），则李代数 \(\mathfrak{s}\mathfrak{u}_{\Delta,\sigma}\) 是 \(\mathrm{A}_l\) 型的绝对单李代数，且相应的同态 \(\pi\) 是与 \(k_{1}\) 相伴的那一个。两个这样的代数 \(\mathfrak{s}\mathfrak{u}_{\Delta,\sigma}\) 与 \(\mathfrak{s}\mathfrak{u}_{\Delta',\sigma'}\) 同构当且仅当存在 \(k\)-同构 \(f: \Delta \to \Delta'\) 使得 \(\sigma' \circ f = f \circ \sigma\)。

证明：当 \(\mathrm{D} = \mathbf{M}_{l+1}(k_{1})\) 时，存在 \(l + 1\) 阶的可逆埃尔米特矩阵 H，在相差一个 \(k^{*}\) 中元素乘积的意义下唯一，使得

\[
\sigma (x) = H. ^ {t} \bar {x}. H ^ {- 1}
\]

对所有 \(x \in \mathbf{M}_{l+1}(k_1)\) 成立；于是代数 g 可与代数 \(\mathfrak{su}(l + 1, H)\)（由满足 \(x.H + H.^t\bar{x} = 0\) 且 \(\operatorname{Tr}(x) = 0\) 的矩阵 x 组成）等同起来。

¶ 17) 令 g 为 \(B_{l}\) 型（分别地，\(C_{l}, D_{l}\) 型）的绝对单李代数，其中 \(l \geq 2\)（分别地，\(l \geq 3, l \geq 4\)）。进一步假设：当 g 为 \(D_{4}\) 型时，§5 习题8 中定义的同态 \(\pi : \operatorname{Gal}(\bar{k}/k) \to \operatorname{Aut}(\bar{\mathrm{R}}, \bar{\mathrm{B}})\) 的像的阶 \(\leq 2\)。证明存在维数为 \((2l + 1)^{2}\)（分别地，\(4l^{2}, 4l^{2}\)）的中心单代数 D 及 D 的一个对合反自同构 \(\sigma\)，使得 g 同构于 D 中由满足 \(\sigma(x) = -x\) 且 \(\operatorname{Tr}_{\mathrm{D}}(x) = 0\) 的元素 x 组成的李子代数（方法同习题16 a)）。\(^{22}\)

\begin{enumerate}
\def\labelenumi{\arabic{enumi})}
\setcounter{enumi}{17}
\tightlist
\item
  令 V 为有限维向量空间，Q 为 V 上的非退化二次型，\(\Psi\) 为与 Q 相伴的对称双线性形式。以 g 表示李代数 \(\mathfrak{o}(\varPsi)\)，并把 \(\varPsi\) 到 \(\bigwedge^{2}(V)\) 上的扩张记作 \(\varPsi_{2}\)。
\end{enumerate}

\begin{enumerate}
\def\labelenumi{\alph{enumi})}
\tightlist
\item
  证明存在向量空间同构 \(\theta : \bigwedge^{2}(\mathrm{V}) \to \mathfrak{g}\)，它由下列等价性质刻画：
\end{enumerate}

\begin{enumerate}
\def\labelenumi{(\roman{enumi})}
\item
  对 \(a, b\) 与 \(x\)（在 \(\mathrm{V}\) 中），有 \(\theta(a \wedge b). x = a. \Psi(x, b) - b. \Psi(x, a)\)。
\item
  对 \(x, y\)（在 V 中）及 \(u\)（在 \(\bigwedge^2 (\mathrm{V})\) 中），有 \(\varPsi_2(x\wedge y,u) = \varPsi (x,\theta (u).y)\)。
\end{enumerate}

令 \(\sigma\) 为 \(\mathfrak{g}\) 在 \(V\) 上的恒等表示；则 \(\theta\) 是 \(\bigwedge^{2}(V)\) 与 \(\mathfrak{g}\) 的伴随表示之间的同构。

\begin{enumerate}
\def\labelenumi{\alph{enumi})}
\setcounter{enumi}{1}
\item
  按第2节引理1 的方式定义线性映射 \(f: \mathfrak{g} \to \mathrm{C}^{+}(\mathrm{Q})\)。证明对 V 中的 a, b 有 \(f\theta(a \wedge b) = \frac{1}{2}(ab - ba)\)，并由此给出第2节引理1 的断言 (iii)、(iv) 与 (v) 的一个新证明。
\item
  记号 \(l, \mathrm{F}, \mathrm{F}', e_0, \mathrm{N}, \lambda\) 与 \(\rho\) 沿用第2节的。取 \(e \neq 0\)（在 \(\bigwedge (\mathrm{F}')\) 中），并用下式定义 \(\varPhi\)（\(\mathrm{N}\) 上的双线性形式）：
\end{enumerate}

\[
x \wedge y = (- 1) ^ {p (p + 1) / 2} \Phi (x, y). e \quad (x \in \bigwedge^ {p} (\mathrm{F} ^ {\prime}), y \in \mathrm{N}).
\]

证明 \(\Phi(\lambda(a).x,y)+\Phi(x,\lambda(a).y)=0\)（对 N 中的 x, y 及 \(F\oplus F'\) 中的 a）。由此推出 \(\Phi\) 在 g 的表示 \(\rho\) 下不变（注意由上面的 a) 与 b)，李代数 \(\rho(\mathfrak{g})\) 由 \(\lambda(F\oplus F')\) 生成）。

\begin{enumerate}
\def\labelenumi{\arabic{enumi})}
\setcounter{enumi}{18}
\tightlist
\item
  令 \(\mathrm{V}\) 为有限维向量空间，\(\mathrm{E} = \mathrm{V} \oplus \mathrm{V}^*\)。用下式定义 \(\varPhi\)（\(\mathrm{E}\) 上的非退化双线性形式）：
\end{enumerate}

\[
\Phi ((x, x ^ {*}), (y, y ^ {*})) = \langle x, y ^ {*} \rangle + \langle y, x ^ {*} \rangle .
\]

置 \(N = \bigwedge(V)\)，并置 \(Q(x) = \frac{1}{2}\Phi(x, x)\)（对 \(x \in E\)）。以 \(\lambda : C(Q) \to \text{End}(N)\) 表示第2节.IV 中的旋量表示，按第2节引理1 定义 \(f : \mathfrak{o}(\Phi) \to C^{+}(Q)\)，并以 \(\rho\) 表示线性表示 \(\lambda \circ f\)（代数 \(\mathfrak{o}(\Phi)\) 在 N 上的）。

\begin{enumerate}
\def\labelenumi{\alph{enumi})}
\item
  对 \(u\)（\(V\) 的任一自同态），用 \(\tilde{u}\)（\(E\) 的自同态）与之对应，公式为 \(\tilde{u}(x, x^{*}) = (u(x), -^{t}u(x^{*}))\)。证明 \(u \mapsto \tilde{u}\) 是从 \(\mathfrak{gl}(V)\) 到 \(\mathfrak{o}(\varPhi)\) 的李代数同态；此外，对 \(u\)（\(\mathfrak{gl}(V)\) 中的），\(\rho(\tilde{u})\) 是代数 \(\bigwedge(V) = N\) 的唯一导子，它与 \(u\)（在 \(V\) 上）一致。
\item
  令 \(\varPsi\) 为 V 上的非退化交错双线性形式，\(\gamma: V \to V^{*}\) 为由 \(\varPsi(x, y) = \langle x, \gamma(y) \rangle\)（x, y 在 V 中）定义的同构。证明由下式定义的 E 的自同态 \(\tilde{X}_{+}\) 与 \(\tilde{X}_{-}\)
\end{enumerate}

\[
\tilde {X} _ {+} (x, x ^ {*}) = (\gamma^ {- 1} (x ^ {*}), 0), \quad \tilde {X} _ {-} (x, x ^ {*}) = (0, - \gamma (x))
\]

属于 \(\mathfrak{o}(\varPhi)\)。置 \(\tilde{H} = (-1)^{\sim}\)。证明 \((\tilde{H}, \tilde{X}_{+}, \tilde{X}_{-})\) 是一个 \(\mathfrak{sl}_2\) 三元组（在李代数 \(\mathfrak{o}(\varPhi)\) 中）。

\begin{enumerate}
\def\labelenumi{\alph{enumi})}
\setcounter{enumi}{2}
\tightlist
\item
  证明 \(\rho\) 把自同态 \(\tilde{H}, \tilde{X}_{+}, \tilde{X}_{-}\)（\(\mathfrak{o}(\Phi)\) 的）分别变为第3节.IV 中记作 \(H, X_{+}\) 与 \(X_{-}\) 的 N 的自同态。由此推出 \((H, X_{+}, X_{-})\) 是一个 \(\mathfrak{sl}_{2}\) 三元组（在李代数 \(\mathfrak{gl}(N)\) 中）。
\end{enumerate}

\specialchapter{半单李代数若干重要性质的综述}

在本综述中，g 表示 k 上的半单李代数。

\section*{嘉当子代数}

\begin{enumerate}
\def\labelenumi{\arabic{enumi})}
\item
  令 E 为 g 中在 g 内为约化的交换子代数所成的集合；它也是 g 中所有元素皆半单的交换子代数所成的集合。g 的嘉当子代数就是 E 中的极大元。
\item
  令 \(x\) 为 \(\mathfrak{g}\) 的正则元。则 \(x\) 是半单元。\(\mathfrak{g}\) 中存在唯一包含 \(x\) 的嘉当子代数；它就是 \(x\) 在 \(\mathfrak{g}\) 中的交换子。
\item
  令 \(x\) 为 \(\mathfrak{g}\) 的半单元。则 \(x\) 属于 \(\mathfrak{g}\) 的某个嘉当子代数。元素 \(x\) 正则当且仅当 \(x\) 的交换子的维数等于 \(\mathfrak{g}\) 的秩。
\item
  令 \(\mathfrak{h}\) 为 \(\mathfrak{g}\) 的嘉当子代数。称 \(\mathfrak{h}\) 是分裂的，若 \(\operatorname{ad}_{\mathfrak{g}}x\) 对所有 \(x \in \mathfrak{h}\) 都可三角化。此外，称 \(\mathfrak{g}\) 是可分裂的，若 \(\mathfrak{g}\) 有分裂嘉当子代数（当 \(k\) 代数闭时即是这种情形）。分裂半单李代数是指一个偶对 \((\mathfrak{g}, \mathfrak{h})\)，其中 \(\mathfrak{g}\) 是半单李代数，\(\mathfrak{h}\) 是 \(\mathfrak{g}\) 的分裂嘉当子代数。
\end{enumerate}

在本综述的其余部分，\((\mathfrak{g},\mathfrak{h})\) 表示分裂半单李代数。

\section*{根系}

\begin{enumerate}
\def\labelenumi{\arabic{enumi})}
\setcounter{enumi}{4}
\item
  对 \(\alpha\)（\(\mathfrak{h}^*\)——\(\mathfrak{h}\) 的对偶——中的任一元素），令 \(\mathfrak{g}^{\alpha}\) 为由 \(x\in \mathfrak{g}\)（满足 \([h,x] = \alpha (h)x\) 对所有 \(h\in \mathfrak{h}\)）所成的集合。若 \(\alpha = 0\)，则 \(\mathfrak{g}^{\alpha} = \mathfrak{h}\)。任何 \(\alpha \in \mathfrak{h}^{*} - \{0\}\)（若它使 \(\mathfrak{g}^{\alpha}\neq 0\)）都称为 \((\mathfrak{g},\mathfrak{h})\) 的根。以 \(\mathrm{R}(\mathfrak{g},\mathfrak{h})\)（或简记为 R）表示 \((\mathfrak{g},\mathfrak{h})\) 的根的集合。它是 \(\mathfrak{h}^*\) 中第VI章§1第4节意义下的既约根系。\(\mathfrak{g}\) 单当且仅当 R 不可约。
\item
  对所有 \(\alpha \in \mathbb{R}\)，\(\mathfrak{g}^{\alpha}\) 的维数为 1。向量空间 \([\mathfrak{g}^{\alpha}, \mathfrak{g}^{-\alpha}]\) 包含于 \(\mathfrak{h}\)，维数为 1，且含有唯一的元素 \(H_{\alpha}\)（满足 \(\alpha(H_{\alpha}) = 2\)）；我们有 \(H_{\alpha} = \alpha^{\vee}\)（第VI章§1第1节）；诸 \(H_{\alpha}\)（对 \(\alpha \in \mathbb{R}\)）所成的集合是逆根系 \(\mathbb{R}^{\vee}\)（\(\mathbb{R}\) 的）。
\item
  我们有 \(\mathfrak{g} = \mathfrak{h} \oplus \bigoplus_{\alpha \in \mathrm{R}} \mathfrak{g}^{\alpha}\)。存在一个族 \((X_{\alpha})_{\alpha \in \mathrm{R}}\)，使得对所有 \(\alpha \in \mathrm{R}\)，有 \(X_{\alpha} \in \mathfrak{g}^{\alpha}\) 且 \([X_{\alpha}, X_{-\alpha}] = -H_{\alpha}\)。每个 \(x \in \mathfrak{g}\) 都可唯一地写成形式
\end{enumerate}

\[
x = h + \sum_ {\alpha \in \mathrm{R}} \lambda_ {\alpha} X _ {\alpha}, \quad \text { where } h \in \mathfrak {h}, \lambda_ {\alpha} \in k.
\]

两个元素的括号可用下列公式计算

\[
\begin{array}{c} [ h, X _ {\alpha} ] = \alpha (h) X _ {\alpha} \\ \relax [ X _ {\alpha}, X _ {\beta} ] = 0 \quad \text {if} \alpha + \beta \notin \mathrm{R} \cup \{0 \} \\ \relax [ X _ {\alpha}, X _ {- \alpha} ] = - H _ {\alpha} \\ \relax [ X _ {\alpha}, X _ {\beta} ] = \mathrm{N} _ {\alpha \beta} X _ {\alpha + \beta} \quad \text {if} \alpha + \beta \in \mathrm{R}, \end{array}
\]

其中 \(\mathrm{N}_{\alpha \beta}\) 是 \(k\) 中的非零元素。

\begin{enumerate}
\def\labelenumi{\arabic{enumi})}
\setcounter{enumi}{7}
\tightlist
\item
  令 B 为 R 的基。代数 g 由诸 \(X_{\alpha}\) 与诸 \(X_{-\alpha}\)（对 \(\alpha \in B\)）生成。我们有 \([X_{\alpha}, X_{-\beta}] = 0\)（若 \(\alpha, \beta \in B\) 且 \(\alpha \neq \beta\)）。令 \((n(\alpha, \beta))_{\alpha, \beta \in B}\) 为 R 的（关于 B 的）嘉当矩阵。我们有 \(n(\alpha, \beta) = \alpha(H_{\beta})\)。若 \(\alpha, \beta \in B\) 且 \(\alpha \neq \beta\)，则 \(n(\alpha, \beta)\) 是负整数，且
\end{enumerate}

\[
(\operatorname{ad} X _ {\beta}) ^ {1 - n (\alpha , \beta)} X _ {\alpha} = 0 \quad \text { and } \quad (\operatorname{ad} X _ {- \beta}) ^ {1 - n (\alpha , \beta)} X _ {- \alpha} = 0.
\]

\begin{enumerate}
\def\labelenumi{\arabic{enumi})}
\setcounter{enumi}{8}
\tightlist
\item
  若 \(\alpha, \beta, \alpha + \beta \in \mathbb{R}\)，令 \(q_{\alpha \beta}\) 为最大整数 \(j\)（使 \(\beta - j\alpha \in \mathbb{R}\)）。7) 中的族 \((X_{\alpha})_{\alpha \in \mathbb{R}}\) 可以取得使 \(\mathrm{N}_{\alpha, \beta} = \mathrm{N}_{-\alpha, -\beta}\)（当 \(\alpha, \beta, \alpha + \beta \in \mathbb{R}\) 时）。此时 \(\mathrm{N}_{\alpha \beta} = \pm (q_{\alpha \beta} + 1)\)。存在一个对合自同构 \(\theta\)（\(\mathfrak{g}\) 的），它把 \(X_{\alpha}\) 变为 \(X_{-\alpha}\)（对所有 \(\alpha \in \mathbb{R}\)）；我们有 \(\theta(h) = -h\)（对所有 \(h \in \mathfrak{h}\)）。\(\mathbf{Z}\)-子模 \(\mathfrak{g}_{\mathbf{Z}}\)（\(\mathfrak{g}\) 中由诸 \(H_{\alpha}\) 与诸 \(X_{\alpha}\) 生成的）是 \(\mathbf{Z}\)-李子代数（\(\mathfrak{g}\) 的），且典则映射 \(\mathfrak{g}_{\mathbf{Z}} \otimes_{\mathbf{Z}} k \to \mathfrak{g}\) 是同构。
\end{enumerate}

偶对 \((\mathfrak{g},\mathfrak{h})\) 可由一个分裂半单 Q-李代数经标量扩张得到。

\begin{enumerate}
\def\labelenumi{\arabic{enumi})}
\setcounter{enumi}{9}
\item
  R 的 Weyl 群、权群、\ldots{}（以及 \ldots{} 等诸群）称为 \((\mathfrak{g},\mathfrak{h})\) 的相应各群。以下把 Weyl 群记作 W。我们考虑它不仅在 \(h^{*}\) 上、而且（通过结构的转移）也在 h 上的作用。若以 \(h_{Q}\)（分别地，\(h_{Q}^{*}\)）表示 h（分别地，\(h^{*}\)）中由诸 \(H_{\alpha}\)（分别地，诸 \(\alpha\)）生成的 Q-向量子空间，则 h（分别地，\(h^{*}\)）可典则地与 \(h_{Q} \otimes_{Q} k\)（分别地，\(h_{Q}^{*} \otimes_{Q} k\)）等同起来，且 \(h_{Q}^{*}\) 可与 \(h_{Q}\) 的对偶等同起来。当我们谈到 R 的 Weyl 房时，理解为它们在 \(h_{R} = h_{Q} \otimes_{Q} R\) 或 \(h_{R}^{*} = h_{Q}^{*} \otimes_{Q} R\) 中。
\item
  令 \(\varPhi\) 为 \(\mathfrak{g}\) 的 Killing 型。若 \(\alpha + \beta \neq 0\)，则 \(\mathfrak{g}^{\alpha}\) 与 \(\mathfrak{g}^{\beta}\) 关于 \(\varPhi\) 正交。\(\varPhi\) 在 \(\mathfrak{g}^{\alpha} \times \mathfrak{g}^{-\alpha}\) 上的限制是非退化的。若 \(x, y \in \mathfrak{h}\)，则 \(\varPhi(x, y) = \sum_{\alpha \in \mathrm{R}} \alpha(x)\alpha(y)\)。我们有 \(\varPhi(H_{\alpha}, H_{\beta}) \in \mathbf{Z}\)。\(\varPhi\) 在 \(\mathfrak{h}\) 上的限制非退化且在 W 下不变；它在 \(\mathfrak{h}_{\mathbf{Q}}\) 上的限制是正定的。
\item
  \((\mathfrak{g},\mathfrak{h})\) 的根系在同构的意义下只依赖于 g 而不依赖于 h。滥用术语，把 Weyl 群、权群、\ldots{}（\((\mathfrak{g},\mathfrak{h})\) 的）也称为 g 的 Weyl 群、权群、\ldots{}。
\end{enumerate}

若 \(R_{1}\) 是既约根系，则存在分裂半单李代数 \((\mathfrak{g}_{1},\mathfrak{h}_{1})\) 使 \(\mathrm{R}(\mathfrak{g}_{1},\mathfrak{h}_{1})\) 同构于 \(R_{1}\)；且它在同构意义下唯一。

于是可分裂半单李代数的分类就化为根系的分类。

\section*{子代数}

\begin{enumerate}
\def\labelenumi{\arabic{enumi})}
\setcounter{enumi}{12}
\item
  若 \(P \subset R\)，置 \(g^{P} = \bigoplus_{\alpha \in P} g^{\alpha}\)，\(h_{P} = \sum_{\alpha \in P} kH_{\alpha}\)。设 \(P \subset R\)，\(h'\) 为 h 的向量子空间，\(a = h' \oplus g^{P}\)。则 a 是 g 的子代数当且仅当 P 是 R 的闭子集且 \(h'\) 含有 \(h_{P \cap (-P)}\)；a 在 g 中约化当且仅当 P = -P；a 可解当且仅当 \(P \cap (-P) = \varnothing\)。
\item
  令 \(\mathbf{P}\) 为 \(\mathbf{R}\) 的闭子集，\(\mathfrak{b} = \mathfrak{h} \oplus \mathfrak{g}^{\mathrm{P}}\)。下列条件等价：
\end{enumerate}

\begin{enumerate}
\def\labelenumi{(\roman{enumi})}
\item
  \(\mathfrak{b}\) 是 \(\mathfrak{g}\) 的极大可解子代数；
\item
  \(\mathrm{P} \cap (-\mathrm{P}) = \varnothing\) 且 \(\mathrm{P} \cup (-\mathrm{P}) = \mathrm{R}\)；
\item
  存在一个房 \(\mathbf{C}\)（属于 \(\mathbf{R}\)）使 \(\mathrm{P} = \mathrm{R}_{+}(\mathrm{C})\)（参见第VI章§1第6节）。
\end{enumerate}

\((\mathfrak{g},\mathfrak{h})\) 的 Borel 子代数是指 g 的一个含 h 且满足上述条件的子代数。g 的子代数 b 若存在 g 的分裂嘉当子代数 \(h'\) 使 b 是 \((\mathfrak{g},\mathfrak{h}')\) 的 Borel 子代数，则称 b 为 g 的 Borel 子代数；若 k 代数闭，这等价于说 b 是 g 的极大可解子代数。

令 \(\mathfrak{b} = \mathfrak{h} \oplus \mathfrak{g}^{\mathrm{R}_{+}(C)}\) 为 \((\mathfrak{g}, \mathfrak{h})\) 的 Borel 子代数。\(\mathfrak{b}\) 的最大幂零理想是 \([\mathfrak{b}, \mathfrak{b}] = \mathfrak{g}^{\mathrm{R}_{+}(C)}\)。令 \(B\) 为 \(R\) 的与 \(C\) 相伴的基；代数 \([\mathfrak{b}, \mathfrak{b}]\) 由诸 \(\mathfrak{g}^{\alpha}\)（\(\alpha \in B\)）生成。

若 \(\mathfrak{b},\mathfrak{b}'\) 是 \(\mathfrak{g}\) 的 Borel 子代数，则存在 \(\mathfrak{g}\) 的嘉当子代数含于 \(\mathfrak{b} \cap \mathfrak{b}'\)；这样的子代数是分裂的。

\begin{enumerate}
\def\labelenumi{\arabic{enumi})}
\setcounter{enumi}{14}
\tightlist
\item
  令 \(\mathbf{P}\) 为 \(\mathbf{R}\) 的闭子集，\(\mathfrak{p} = \mathfrak{h} \oplus \mathfrak{g}^{\mathrm{P}}\)。下列条件等价：
\end{enumerate}

\begin{enumerate}
\def\labelenumi{(\roman{enumi})}
\item
  p 含有 \((\mathfrak{g}, \mathfrak{h})\) 的一个 Borel 子代数；
\item
  \(P \cup (-P) = R;\)
\item
  存在 R 的房 C 使 \(\mathrm{P} \supset \mathrm{R}_{+}(\mathrm{C})\)。
\end{enumerate}

\((\mathfrak{g},\mathfrak{h})\) 的抛物子代数是指 g 的一个含 h 且满足上述条件的子代数。g 的子代数 p 若存在 g 的分裂嘉当子代数 \(h'\) 使 p 是 \((\mathfrak{g},\mathfrak{h}')\) 的抛物子代数，则称 p 是抛物的。

令 \(\mathfrak{p} = \mathfrak{h} \oplus \mathfrak{g}^{\mathrm{P}}\) 为 \((\mathfrak{g}, \mathfrak{h})\) 的抛物子代数，Q 为由 \(\alpha \in \mathrm{P}\)（满足 \(-\alpha \notin \mathrm{P}\)）组成的集合，\(\mathfrak{s} = \mathfrak{h} \oplus \mathfrak{g}^{\mathrm{P} \cap (-\mathrm{P})}\)。则 \(\mathfrak{p} = \mathfrak{s} \oplus \mathfrak{g}^{\mathrm{Q}}\)，\(\mathfrak{s}\) 在 \(\mathfrak{g}\) 中约化，且 \(\mathfrak{g}^{\mathrm{Q}}\) 是 \(\mathfrak{p}\) 的最大幂零理想兼 \(\mathfrak{p}\) 的幂零根基。\(\mathfrak{p}\) 的中心为 0。

\section*{自同构}

\begin{enumerate}
\def\labelenumi{\arabic{enumi})}
\setcounter{enumi}{15}
\tightlist
\item
  \(\mathrm{Aut}(\mathfrak{g})\) 中由诸 \(e^{\mathrm{ad}x}\)（\(x\) 幂零）生成的子群是初等自同构群 \(\mathrm{Aut}_e(\mathfrak{g})\)（\(\mathfrak{g}\) 的）；它是 \(\mathrm{Aut}(\mathfrak{g})\) 的正规子群；它等于自身的换位子群。
\end{enumerate}

若 \(\bar{k}\) 是 \(k\) 的代数闭包，则群 \(\operatorname{Aut}(\mathfrak{g})\) 自然地嵌入 \(\operatorname{Aut}(\mathfrak{g} \otimes_k \bar{k})\)。置

\[
\operatorname{Aut} _ {0} (\mathfrak {g}) = \operatorname{Aut} (\mathfrak {g}) \cap \operatorname{Aut} _ {e} (\mathfrak {g} \otimes_ {k} \bar {k});
\]

这是 \(\operatorname{Aut}(\mathfrak{g})\) 的正规子群，与 \(\bar{k}\) 的选取无关。我们有

\[
\operatorname{Aut} _ {e} (\mathfrak {g}) \subset \operatorname{Aut} _ {0} (\mathfrak {g}) \subset \operatorname{Aut} (\mathfrak {g}).
\]

\(\mathrm{Aut}_0(\mathfrak{g})\) 的换位子群是 \(\mathrm{Aut}_e(\mathfrak{g})\)。在 Zariski 拓扑下，\(\mathrm{Aut}(\mathfrak{g})\) 与 \(\mathrm{Aut}_0(\mathfrak{g})\) 在 \(\mathrm{End}_k(\mathfrak{g})\) 中闭，\(\mathrm{Aut}_0(\mathfrak{g})\) 是 \(\mathrm{Aut}(\mathfrak{g})\) 的单位元连通分支，且 \(\mathrm{Aut}_e(\mathfrak{g})\) 在 \(\mathrm{Aut}_0(\mathfrak{g})\) 中稠密。

令 B 为 R 的基，\(\operatorname{Aut}(R,B)\) 为使 B 稳定的 R 的自同构群。则 \(\operatorname{Aut}(\mathfrak{g})\) 是一个同构于 \(\operatorname{Aut}(R,B)\) 的子群与 \(\operatorname{Aut}_{0}(\mathfrak{g})\) 的半直积；特别地，\(\operatorname{Aut}(\mathfrak{g})/\operatorname{Aut}_{0}(\mathfrak{g})\) 同构于 \(\operatorname{Aut}(R,B)\)，而后者本身又同构于 g 的邓金图的一个自同构群。

\begin{enumerate}
\def\labelenumi{\arabic{enumi})}
\setcounter{enumi}{16}
\tightlist
\item
  g 的标架是指一个三元组 \((\mathfrak{h}^{\prime},\mathrm{B},(X_{\alpha})_{\alpha\in\mathrm{B}})\)，其中 \(h^{\prime}\) 是 g 的分裂嘉当子代数，B 是 \(\mathrm{R}(\mathfrak{g},\mathfrak{h}^{\prime})\) 的基，且对所有 \(\alpha\in B\)，\(X_{\alpha}\) 是 \(g^{\alpha}\) 的非零元素。群 \(\operatorname{Aut}_{0}(\mathfrak{g})\) 单迁移地作用于 g 的标架所成的集合上。
\end{enumerate}

群 \(\mathrm{Aut}_e(\mathfrak{g})\) 迁移地作用于偶对 \((\mathfrak{k},\mathfrak{b})\) 所成的集合上，其中 \(\mathfrak{k}\) 是 \(\mathfrak{g}\) 的分裂嘉当子代数，\(\mathfrak{b}\) 是 \((\mathfrak{g},\mathfrak{k})\) 的 Borel 子代数。

\begin{enumerate}
\def\labelenumi{\arabic{enumi})}
\setcounter{enumi}{17}
\tightlist
\item
  以 \(\operatorname{Aut}(\mathfrak{g},\mathfrak{h})\) 表示由 \(s\in \operatorname{Aut}(\mathfrak{g})\)（满足 \(s(\mathfrak{h}) = \mathfrak{h}\)）组成的集合。置
\end{enumerate}

\[
\operatorname{Aut} (\mathfrak {g}, \mathfrak {h}) = \operatorname{Aut} _ {e} (\mathfrak {g}) \cap \operatorname{Aut} (\mathfrak {g}, \mathfrak {h}), \quad \operatorname{Aut} _ {0} (\mathfrak {g}, \mathfrak {h}) = \operatorname{Aut} _ {0} (\mathfrak {g}) \cap \operatorname{Aut} (\mathfrak {g}, \mathfrak {h}).
\]

若 \(s \in \text{Aut}(\mathfrak{g}, \mathfrak{h})\)，则 \(s | \mathfrak{h}\) 的反步映射是 R 的自同构群 \(\text{A(R)}\) 的一个元素；以 \(\varepsilon(s)\) 记这个元素；映射 \(\varepsilon\) 是从 \(\text{Aut}(\mathfrak{g}, \mathfrak{h})\) 到 \(\text{A(R)}\) 的同态。我们有 \(\text{Aut}_{0}(\mathfrak{g}) = \text{Aut}_{e}(\mathfrak{g})\). Ker \(\varepsilon\)，且

\[
\varepsilon (\operatorname{Aut} _ {0} (\mathfrak {g}, \mathfrak {h})) = \varepsilon (\operatorname{Aut} _ {e} (\mathfrak {g}, \mathfrak {h})) = W.
\]

令 \(\mathrm{T}_{\mathrm{P}} = \mathrm{Hom}(\mathrm{P}(\mathrm{R}), k^{*})\)，\(\mathrm{T}_{\mathrm{Q}} = \mathrm{Hom}(\mathrm{Q}(\mathrm{R}), k^{*})\)。\(\mathrm{Q}(\mathrm{R})\) 到 \(\mathrm{P}(\mathrm{R})\) 的单射定义了一个从 \(\mathrm{T}_{\mathrm{P}}\) 到 \(\mathrm{T}_{\mathrm{Q}}\) 的同态；以 \(\mathrm{Im}(\mathrm{T}_{\mathrm{P}})\) 表示其像。若 \(t \in \mathrm{T}_{\mathrm{Q}}\)，令 \(f(t)\) 为 \(\mathfrak{g}\) 的满足下述条件的自同态：对所有 \(\alpha \in \mathrm{R} \cup \{0\}\)，\(f(t)|\mathfrak{g}^{\alpha}\) 是比为 \(t(\alpha)\) 的位似；我们有 \(f(t)\in \mathrm{Aut}_{0}(\mathfrak{g},\mathfrak{h})\)，且 \(f\) 是从 \(\mathrm{T}_{\mathrm{Q}}\) 到 \(\mathrm{Aut}_0(\mathfrak{g},\mathfrak{h})\) 的单同态。序列

\[
\{1 \} \longrightarrow \mathrm {T_ {Q}} \stackrel {{f}} {{\longrightarrow}} \operatorname{Aut} (\mathfrak {g}, \mathfrak {h}) \stackrel {{\varepsilon}} {{\longrightarrow}} \mathrm{A(R)} \longrightarrow \{1 \}
\]

与

\[
\{1 \} \longrightarrow \mathrm{T} _ {\mathrm{Q}} \stackrel {f} {\longrightarrow} \mathrm{Aut} _ {0} (\mathfrak {g}, \mathfrak {h}) \stackrel {\varepsilon} {\longrightarrow} \mathrm{W} \longrightarrow \{1 \}
\]

是正合的。我们有 \(f(\mathrm{Im}(\mathrm{T}_{\mathrm{P}})) \subset \mathrm{Aut}_{e}(\mathfrak{g}, \mathfrak{h})\)；f 通过过渡到商定义出一个满射\(^{23}\)同态 \(\mathrm{T}_{\mathrm{Q}}/\mathrm{Im}(\mathrm{T}_{\mathrm{P}}) \to \mathrm{Aut}_{0}(\mathfrak{g})/\mathrm{Aut}_{e}(\mathfrak{g})\)。在 Zariski 拓扑下，\(f(\mathrm{T}_{\mathrm{Q}})\) 在 \(\mathrm{Aut}(\mathfrak{g})\) 中闭，且 \(f(\mathrm{Im}(\mathrm{T}_{\mathrm{P}}))\) 在 \(f(\mathrm{T}_{\mathrm{Q}})\) 中稠密。

\section*{有限维模}

\begin{enumerate}
\def\labelenumi{\arabic{enumi})}
\setcounter{enumi}{18}
\item
  令 V 为有限维 g-模。对所有 \(\mu\in h^{*}\)，令 \(V^{\mu}\) 为由 \(v\in V\)（满足 \(h.v=\mu(h)v\) 对所有 \(h\in h\)）组成的集合。\(V^{\mu}\) 的维数称为 \(\mu\) 在 V 中的重数；若它 \(\geq1\)，即若 \(V^{\mu}\neq0\)，则称 \(\mu\) 是 V 的权。我们有 \(V=\bigoplus_{\mu\in h^{*}}V^{\mu}\)。V 的每个权都属于 P(R)。若 \(\mu\) 是 V 的权且 \(w\in W\)，则 \(w\mu\) 也是 V 的权，且重数与 \(\mu\) 相同。若 \(v\in V^{\mu}\) 且 \(x\in g^{\alpha}\)，则 \(x.v\in V^{\mu+\alpha}\)。
\item
  令 B 为 R 的基。给定 B 就确定了 \(h_{Q}^{*}\) 上的一个序关系：\(h_{Q}^{*}\) 中 \(\geq 0\) 的元素是 B 中元素以 \(\geq 0\) 的有理数为系数的线性组合。以 \(Q_{+}(R)\)（分别地，\(R_{+}\)）表示 \(Q(R)\)（分别地，R）的正元素所成的集合。
\end{enumerate}

令 \(V\) 为有限维单 \(\mathfrak{g}\)-模。则 \(V\) 有最高权 \(\lambda\)。这个权的重数为 1，且它是支配权：若 \(\alpha \in R_{+}\)，\(\lambda(H_{\alpha})\) 是 \(\geq 0\) 的整数。我们有 \(\mathfrak{g}^{\alpha}V^{\lambda} = 0\)（若 \(\alpha \in R_{+}\)）。\(V\) 的每个权都具有 \(\lambda - \nu\)（\(\nu \in Q_{+}(R)\)）的形式；反之，若一个权是支配权且具有 \(\lambda - \nu\)（\(\nu \in Q_{+}(R)\)）的形式，则它是 \(V\) 的权。

\begin{enumerate}
\def\labelenumi{\arabic{enumi})}
\setcounter{enumi}{20}
\tightlist
\item
  具有相同最高权的两个有限维单 g-模同构。每个支配权都是某个有限维单 g-模的最高权。
\end{enumerate}

每个有限维单 g-模都是绝对单的。

\begin{enumerate}
\def\labelenumi{\arabic{enumi})}
\setcounter{enumi}{21}
\item
  令 \(\varPhi\) 为 g 的 Killing 型，\(C \in U(\mathfrak{g})\) 为相应的 Casimir 元，\(\langle\cdot,\cdot\rangle\) 为 \(h^{*}\) 上的、\(\varPhi|h \times h\) 的逆形式，\(\rho = \frac{1}{2} \sum_{\alpha \in R_{+}} \alpha\)。令 V 为有限维单 g-模，最高权为 \(\lambda\)。则 \(C_{V}\) 是比为 \(\langle\lambda, \lambda + 2\rho\rangle\) 的位似。
\item
  令 \(\mathrm{V}\) 为有限维 \(\mathfrak{g}\)-模，\(\mathrm{V}^*\) 为其对偶。则 \(\mu \in \mathfrak{h}^*\) 是 \(\mathrm{V}^*\) 的权当且仅当 \(-\mu\) 是 \(\mathrm{V}\) 的权，且 \(\mu \in \mathrm{V}^*\) 的重数等于 \(-\mu\) 在 V 中的重数。若 V 是最高权为 \(\lambda\) 的单模，则 \(V^{*}\) 是最高权为 \(-w_{0}\lambda\) 的单模，其中 \(w_{0}\) 是 W 中把 B 变为 -B 的元素。
\item
  令 \(\mathrm{V}\) 为有限维单 \(\mathfrak{g}\)-模，其最高权为 \(\lambda\)，\(\mathcal{B}\) 为 \(\mathfrak{g}\)-不变双线性形式（\(\mathrm{V}\) 上的）所成的向量空间。令 \(m\) 为整数 \(\sum_{\alpha \in \mathbb{R}_+} \lambda(H_\alpha)\)，并令 \(w_0 \in \mathrm{W}\) 如 23) 中那样。若 \(w_0 \lambda \neq -\lambda\)，则 \(\mathrm{V}\) 与 \(V^*\) 不同构，且 \(\mathcal{B} = 0\)。若 \(w_0 \lambda = -\lambda\)，则 \(\dim \mathcal{B} = 1\)，且 \(\mathcal{B}\) 的每个非零元素都非退化；若 \(m\) 为偶（分别地，奇）数，则 \(\mathcal{B}\) 的每个元素都是对称（分别地，交错）的。
\item
  令 \(\mathbf{Z}[\mathrm{P}]\) 为群 \(\mathrm{P} = \mathrm{P}(\mathrm{R})\) 的系数在 \(\mathbf{Z}\) 中的群代数。若 \(\lambda \in \mathrm{P}\)，以 \(e^{\lambda}\) 表示 \(\mathbf{Z}[\mathrm{P}]\) 中相应的元素；诸 \(e^{\lambda}\)（\(\lambda \in \mathrm{P}\)）构成一个 \(\mathbf{Z}\)-基（\(\mathbf{Z}[\mathrm{P}]\) 的），且 \(e^{\lambda}e^{\mu} = e^{\lambda +\mu}\)（对 \(\lambda, \mu \in \mathrm{P}\)）。
\end{enumerate}

令 \(\mathrm{V}\) 为有限维 \(\mathfrak{g}\)-模。\(\mathrm{V}\) 的特征标记作 \(\operatorname{ch} \mathrm{V}\)，它是元素 \(\sum_{\mu \in \mathrm{P}} (\dim \mathrm{V}^{\mu}) e^{\mu}\)（\(\mathbf{Z}[\mathrm{P}]\) 的）；这个元素属于子代数 \(\mathbf{Z}[\mathrm{P}]^{\mathrm{W}}\)——\(\mathbf{Z}[\mathrm{P}]\) 中由在 \(\mathrm{W}\) 下不变的元素组成的。我们有

\[
\operatorname{ch} \left(\mathrm{V} \oplus \mathrm{V} ^ {\prime}\right) = \operatorname{ch} \mathrm{V} + \operatorname{ch} \mathrm{V} ^ {\prime} \quad \text { and } \quad \operatorname{ch} \left(\mathrm{V} \otimes \mathrm{V} ^ {\prime}\right) = (\operatorname{ch} \mathrm{V}). (\operatorname{ch} \mathrm{V} ^ {\prime}).
\]

具有相同特征标的两个有限维 g-模同构。

对所有 \(\alpha\in B\)，令 \(V_{\alpha}\) 为以基本权 \(\varpi_{\alpha}\)（对应于 \(\alpha\)）为最高权的单 g-模。诸元素 ch \(V_{\alpha},\alpha\in B\) 代数无关，且生成 Z-代数 \(Z[P]^{W}\)。

\begin{enumerate}
\def\labelenumi{\arabic{enumi})}
\setcounter{enumi}{25}
\tightlist
\item
  令 \(\rho\) 为诸根 \(\geq 0\) 的根之和的一半。对所有 \(w \in W\)，令 \(\varepsilon(w)\) 为 w 的行列式，等于 \(\pm 1\)。若 V 是最高权为 \(\lambda\) 的有限维单 g-模，则
\end{enumerate}

\[
\left(\sum_ {w \in \mathrm{W}} \varepsilon (w) e ^ {w \rho}\right). \mathrm{chV} = \sum_ {w \in \mathrm{W}} \varepsilon (w) e ^ {w (\lambda + \rho)},
\]

以及

\[
\dim V = \prod_ {\alpha \in R _ {+}} \frac {\langle \lambda + \rho , H _ {\alpha} \rangle}{\langle \rho , H _ {\alpha} \rangle}.
\]

\begin{enumerate}
\def\labelenumi{\arabic{enumi})}
\setcounter{enumi}{26}
\tightlist
\item
  对所有 \(\nu \in \mathrm{P}\)，令 \(\mathfrak{P}(\nu)\) 为族 \((n_{\alpha})_{\alpha \in \mathrm{R}_{+}}\) 的个数，其中诸 \(n_{\alpha}\) 是 \(\geq 0\) 的整数且满足 \(\nu = \sum_{\alpha \in \mathrm{R}_{+}} n_{\alpha} \alpha\)。令 \(\mathrm{V}\) 为有限维单 \(\mathfrak{g}\)-模，其最高权为 \(\lambda\)。若 \(\mu \in \mathrm{P}\)，则 \(\mu\) 在 \(\mathrm{V}\) 中的重数是
\end{enumerate}

\[
\sum_ {w \in \mathrm{W}} \varepsilon (w) \mathfrak {P} (w (\lambda + \rho) - (\mu + \rho)).
\]

\begin{enumerate}
\def\labelenumi{\arabic{enumi})}
\setcounter{enumi}{27}
\tightlist
\item
  令 \(\mathrm{V}, \mathrm{V}', \mathrm{V}''\) 为有限维单 \(\mathfrak{g}\)-模，\(\lambda, \mu, \nu\) 为它们的最高权。在 \(\mathrm{V} \otimes \mathrm{V}'\) 中，\(\mathrm{V}''\) 型的同构型分支的长度为
\end{enumerate}

\[
\sum_ {w, w ^ {\prime} \in \mathrm{W}} \varepsilon (w w ^ {\prime}) \mathfrak {P} (w (\lambda + \rho) + w ^ {\prime} (\mu + \rho) - (\nu + 2 \rho)).
\]

特别地，若 \(\nu = \lambda + \mu\)，则所讨论的同构型分支是单的，且由 \((\mathrm{V} \otimes \mathrm{V}')^{\lambda + \mu} = \mathrm{V}^{\lambda} \otimes \mathrm{V}'^{\mu}\) 生成。

\section*{不变多项式函数}

\begin{enumerate}
\def\labelenumi{\arabic{enumi})}
\setcounter{enumi}{28}
\item
  \(\mathfrak{g}\) 上的多项式函数代数可与对称代数 \(\mathbf{S}(\mathfrak{g}^*)\)（\(\mathfrak{g}^*\) 的）等同起来，从而以典则方式成为一个 \(\mathfrak{g}\)-模；于是有 \(\mathfrak{g}\) 上不变多项式函数的概念。令 \(f \in \mathbf{S}(\mathfrak{g}^*)\)。则 \(f\) 不变当且仅当 \(f \circ s = f\)（对所有 \(s \in \mathrm{Aut}_0(\mathfrak{g})\)），或 \(f \circ s = f\)（对所有 \(s \in \mathrm{Aut}_e(\mathfrak{g})\)）。
\item
  令 \(\mathrm{I}(\mathfrak{g}^*)\) 为 \(\mathfrak{g}\) 上的不变多项式函数代数，\(\mathbf{S}(\mathfrak{h}^*)^{\mathrm{W}}\) 为 \(\mathfrak{h}\) 上的 W-不变多项式函数代数。令
\end{enumerate}

\[
i: \mathbf {S} (\mathfrak {g} ^ {*}) \to \mathbf {S} (\mathfrak {h} ^ {*})
\]

为限制同态。映射 \(i|I(\mathfrak{g}^{*})\) 是从 \(I(\mathfrak{g}^{*})\) 到 \(\mathbf{S}(\mathfrak{h}^{*})^{\mathrm{W}}\) 的同构。若 l 是 g 的秩，则存在 \(I(\mathfrak{g}^{*})\) 的 l 个代数无关的齐次元素，它们生成代数 \(I(\mathfrak{g}^{*})\)。

\begin{enumerate}
\def\labelenumi{\arabic{enumi})}
\setcounter{enumi}{30}
\item
  元素 \(a\)（\(\mathfrak{g}\) 的）幂零当且仅当 \(f(a) = 0\)（对每个齐次元素 \(f\)（属 \(\mathrm{I}(\mathfrak{g}^*)\)，次数 \(> 0\)））。
\item
  令 \(s \in \mathrm{Aut}(\mathfrak{g})\)。则 \(s\) 属于 \(\mathrm{Aut}_0(\mathfrak{g})\) 当且仅当 \(f \circ s = f\)（对所有 \(f \in \mathrm{I}(\mathfrak{g}^*)\)）。
\end{enumerate}

\section*{\texorpdfstring{\(\mathfrak{sl}_2\) 三元组}{\textbackslash mathfrak\{sl\}\_2 三元组}}

\begin{enumerate}
\def\labelenumi{\arabic{enumi})}
\setcounter{enumi}{32}
\item
  g 中的 \(sl_{2}\) 三元组是指一个元素序列 \((x,h,y)\)——g 中异于 \((0,0,0)\) 且满足 \([h,x]=2x,[h,y]=-2y,[x,y]=-h\) 的。此时 x,y 在 g 中幂零，h 在 g 中半单。
\item
  令 \(x\) 为 \(\mathfrak{g}\) 的非零幂零元。存在 \(h, y \in \mathfrak{g}\) 使 \((x, h, y)\) 是一个 \(\mathfrak{sl}_2\) 三元组。
\item
  令 \((x,h,y)\) 与 \((x',h',y')\) 为两个 \(\mathfrak{sl}_2\) 三元组（在 \(\mathfrak{g}\) 中）。下列条件等价：
\end{enumerate}

\begin{enumerate}
\def\labelenumi{\alph{enumi})}
\item
  存在 \(s \in \text{Aut}_{e}(\mathfrak{g})\) 使 \(sx = x'\)；
\item
  存在 \(s \in \mathrm{Aut}_e(\mathfrak{g})\) 使 \(sx = x', sh = h', sy = y'\)。
\end{enumerate}

\begin{enumerate}
\def\labelenumi{\arabic{enumi})}
\setcounter{enumi}{35}
\tightlist
\item
  若 \(k\) 代数闭，则 35) 的条件 \(a)\) 与 \(b)\) 等价于：\(c)\) 存在 \(s \in \operatorname{Aut}_e(\mathfrak{g})\) 使 \(sh = h'\)。
\end{enumerate}

此外，g 的非零幂零元关于 \(\operatorname{Aut}_{e}(\mathfrak{g})\) 的共轭类个数至多等于 \(3^{l}\)，其中 l 是 g 的秩。

\begin{enumerate}
\def\labelenumi{\arabic{enumi})}
\setcounter{enumi}{36}
\tightlist
\item
  g 的幂零元 x 若其中心化子的维数等于 g 的秩，则称为主幂零元。g 中存在主幂零元。若 k 代数闭，则 g 的主幂零元在 \(\operatorname{Aut}_{e}(\mathfrak{g})\) 下共轭。
\end{enumerate}

\chapter{紧实李群}

在本章\(^{1}\)中，「李群」一词指「实数域上的有限维李群」，除另有说明外，「李代数」一词指「实数域上的有限维李代数」，「实李代数」（分别地，「复李代数」）指「实数域上的有限维李代数」（分别地，「复数域上的有限维李代数」）。

我们以 \(G_{0}\) 表示拓扑群 G 的单位元连通分支。以 \(\mathrm{C(G)}\) 表示群 G 的中心，以 \(\mathrm{D(G)}\) 表示其换位子群，以 \(\mathrm{N_{G}(H)}\) 或 \(\mathrm{N(H)}\)（分别地，\(\mathrm{Z_{G}(H)}\) 或 \(\mathrm{Z(H)}\)）表示群 G 的子集 H 的正规化子（分别地，中心化子）。

\section{紧李代数}

\subsection{不变埃尔米特型}

在本节中，字母 k 表示域 R 或 C。令 V 为有限维 k-向量空间，\(\Phi\) 为 V 上分离的\(^{2}\)正埃尔米特型，G 为群，g 为 R-李代数，\(\rho: G \to \mathbf{GL}(V)\) 为群同态，\(\varphi: \mathfrak{g} \to \mathfrak{gl}(V)\) 为 R-李代数同态。

\begin{enumerate}
\def\labelenumi{\alph{enumi})}
\tightlist
\item
  形式 \(\Phi\) 在 G（分别地，g）下不变当且仅当 \(\rho(g)\) 关于 \(\Phi\) 是酉的（对所有 \(g \in G\)）（分别地，\(\varphi(x)\) 是反埃尔米特的\(^{3}\)（关于 \(\Phi\)，对所有 \(x \in g\)））。事实上，以 \(a^{*}\) 表示 V 的自同态 a 关于 \(\Phi\) 的伴随；对 G 中的 g、g 中的 x、V 中的 u 与 v，我们有
\end{enumerate}

\[
\begin{array}{c} \Phi (\rho (g) u, \rho (g) v) = \Phi (\rho (g) ^ {*} \rho (g) u, v), \\ \Phi (\varphi (x) u, v) + \Phi (u, \varphi (x) v) = \Phi ((\varphi (x) + \varphi (x) ^ {*}). u, v); \end{array}
\]

于是，\(\Phi(\rho(g)u,\rho(g)v)=\Phi(u,v)\)（对所有 \(u,v\)（在 V 中））当且仅当 \(\rho(g)^{*}\rho(g)=\mathrm{Id}_{\mathrm{V}}\)；类似地，\(\Phi(\varphi(x)u,v)+\Phi(u,\varphi(x)v)=0\)（对所有 \(u,v\)（在 V 中））当且仅当 \(\varphi(x)+\varphi(x)^{*}=0\)，由此即得所述断言。

\begin{enumerate}
\def\labelenumi{\alph{enumi})}
\setcounter{enumi}{1}
\item
  若形式 \(\varPhi\) 在 G（分别地，g）下不变，则 V 的稳定子空间的正交补是稳定的；特别地，表示 \(\rho\)（分别地，\(\varphi\)）此时是半单的（参见《代数学》第IX章）；此外，对所有 \(g \in G\)（分别地，\(x \in g\)），V 的自同态 \(\rho(g)\)（分别地，\(\varphi(x)\)）此时是半单的，其特征值的绝对值为 1（分别地，特征值都是纯虚数）；事实上，\(\rho(g)\) 是酉的（分别地，\(i\varphi(x)\) 是埃尔米特的，参见《代数学》第IX章）。
\item
  假设 \(k = \mathbf{R}\)。若 \(G\) 是连通李群，\(\rho\) 是李群态射，\(\mathfrak{g}\) 是 \(G\) 的李代数，\(\varphi\) 是由 \(\rho\) 诱导的同态，则 \(\varPhi\) 在 \(G\) 下不变当且仅当它在 \(\mathfrak{g}\) 下不变（第III章§6第5节推论3）。
\item
  V 上存在在 G 下不变的分离的正埃尔米特形式，当且仅当子群 \(\rho(\mathrm{G})\)（\(\mathbf{GL}(\mathrm{V})\) 的）是相对紧的（《积分学》第VII章§3第1节命题1）。
\end{enumerate}

\subsection{连通交换实李群}

令 \(G\) 为连通交换（实）李群。指数映射

\[
\exp_ {\mathrm{G}}: \mathrm{L} (\mathrm{G}) \to \mathrm{G}
\]

是李群态射，且是满射、有离散核（第III章§6第4节命题11），由此可知 \(\mathrm{L}(\mathrm{G})\) 是 \(\mathrm{G}\) 的连通覆盖。

\begin{enumerate}
\def\labelenumi{\alph{enumi})}
\item
  下列条件等价：G 单连通，\(\exp_{G}\) 是同构，G 同构于 \(R^{n}\)（\(n = \dim G\)）。在这种情形，用同构 \(\exp_{G}\) 把 L(G) 的向量空间结构转移到 G 上，就得到 G 上的一个向量空间结构，它是唯一与 G 的拓扑群结构相容的向量空间结构。单连通交换李群称为向量（李）群；除另有说明外，它们总被赋予上面定义的 R-向量空间结构。
\item
  以 \(\Gamma(\mathrm{G})\) 表示 \(\exp_{G}\) 的核。由《一般拓扑学》第VII章§1第1节定理1，群 G 紧当且仅当 \(\Gamma(\mathrm{G})\) 是 \(\mathrm{L}(\mathrm{G})\) 中的格，换句话说（同处所引），当且仅当自由 Z-模 \(\Gamma(\mathrm{G})\) 的秩等于 G 的维数。反之，若 L 是有限维 R-向量空间，\(\Gamma\) 是 L 中的格，则商拓扑群 \(L/\Gamma\) 是紧连通交换李群。
\end{enumerate}

紧连通交换李群称为实环面，或（在本章中）环面。

\begin{enumerate}
\def\labelenumi{\alph{enumi})}
\setcounter{enumi}{2}
\tightlist
\item
  在一般情形，令 E 为 L(G) 中由 \(\Gamma(\mathrm{G})\) 生成的向量子空间，令 V 为其补子空间。则 G 是它的李子群 \(\exp(\mathrm{E})\) 与 \(\exp(\mathrm{V})\) 的直积；前者是环面，后者是向量群。最后，G 的每个紧子群都含于 \(\exp(\mathrm{E})\)（因为它在 \(\exp(V)\) 上的投影必然化为单位元）；于是，子群 \(\exp(E)\) 是 G 的唯一的极大紧子群。
\end{enumerate}

例如，取 \(\mathrm{G} = \mathrm{C}^*\)；把 \(\mathrm{L(G)}\) 与 \(\mathbf{C}\) 等同起来，使 \(\mathrm{G}\) 的指数映射为 \(x \mapsto e^x\)。则 \(\Gamma(\mathrm{G}) = 2\pi i\mathbf{Z}, \mathrm{E} = i\mathbf{R}\)，从而 \(\exp(\mathrm{E}) = \mathbf{U}\)；若取 \(\mathrm{V} = \mathbf{R}\)，则 \(\exp(\mathrm{V}) = \mathbf{R}_+^*\)，我们重新得到《一般拓扑学》第VIII章§1第3节命题1 中构造的同构 \(\mathbf{C}^* \to \mathbf{U} \times \mathbf{R}_+^*\)。

\(d\) ) 最后注意，\(\exp_{\mathrm{G}}:\mathrm{L}(\mathrm{G})\to \mathrm{G}\) 是 \(\mathbf{G}\) 的泛覆盖，因而 \(\Gamma (\mathbf{G})\) 可自然地与 \(\mathbf{G}\) 的基本群等同起来。

\subsection{紧李代数}

命题 1. 令 \(\mathfrak{g}\) 为（实）李代数。下列条件等价：

\begin{enumerate}
\def\labelenumi{(\roman{enumi})}
\item
  g 同构于某个紧李群的李代数。
\item
  群 \(\operatorname{Int}(\mathfrak{g})\)（第III章§6第2节定义2）是紧的。
\item
  \(\mathfrak{g}\) 有一个对称、正且分离的不变双线性形式（第I章§3第6节）。
\item
  g 是约化的（第I章§6第4节定义4）；对所有 \(x \in g\)，自同态 ad x 半单，且特征值都是纯虚数。
\item
  \(\mathfrak{g}\) 约化且其 Killing 型 B 是负的。
\item
  \(\Longrightarrow\) (ii)：若 g 是紧李群 G 的李代数，则群 \(\operatorname{Int}(\mathfrak{g})\) 是分离的，且同构于紧群 \(G_{0}\) 的一个商（第III章§6第4节推论4），因而是紧的。
\item
  \(\Longrightarrow\) (iii)：若群 \(\operatorname{Int}(\mathfrak{g})\) 紧，则 g 上存在一个对称双线性形式，它正、分离且在 \(\operatorname{Int}(\mathfrak{g})\) 下不变（第1节），从而也在 g 的伴随表示下不变。
\item
  \(\Longrightarrow\) (iv)：若 (iii) 成立，则 g 的伴随表示半单（第1节），因而 g 约化；此外，\(x \in g\) 的自同态 ad x 具有所述性质（第1节）。
\item
  \(\Longrightarrow\) (v)：对所有 \(x\in \mathfrak{g}\)，\(\mathrm{B}(x,x) = \mathrm{Tr}((\mathrm{ad}x)^2)\)；于是 \(\mathrm{B}(x,x)\) 是 ad \(x\) 的特征值的平方和，因而当这些特征值都是纯虚数时它是负的。
\item
  \(\Longrightarrow\) (i)：假设 g 约化，从而是一个交换子代数 c 与一个半单子代数 s 的积（第I章§6第4节命题5）。s 的 Killing 型是形式 B 在 s 上的限制，因而当 B 负时它是负且分离的。子群 \(\operatorname{Int}(\mathfrak{s})\)（\(\mathbf{GL}(\mathfrak{s})\) 的）是闭的（它是 \(\operatorname{Aut}(\mathfrak{s})\) 的单位元连通分支，第III章§10第2节推论2）且使分离的正形式 -B 不变；于是它是紧的，且 s 同构于紧李群 \(\operatorname{Int}(\mathfrak{s})\) 的李代数。再者，由于 c 交换，它同构于某个环面 T 的李代数。于是 g 同构于紧李群 \(\operatorname{Int}(\mathfrak{s}) \times \operatorname{T}\) 的李代数。
\end{enumerate}

定义 1. 紧李代数\(^{4}\)是指具有命题 1 中性质 (i) 至 (v) 的李代数。

于是，紧李代数就是交换代数与紧半单代数的积。换句话说，李代数紧当且仅当它约化且其导李代数紧。

紧李群的李代数是紧的。

命题 2. a) 有限个李代数的积是紧李代数，当且仅当每个因子都紧。

\begin{enumerate}
\def\labelenumi{\alph{enumi})}
\setcounter{enumi}{1}
\item
  紧李代数的子代数是紧的。
\item
  令 \(\mathfrak{h}\) 为紧李代数 \(\mathfrak{g}\) 的理想。则代数 \(\mathfrak{g} / \mathfrak{h}\) 紧，且扩张 \(\mathfrak{h} \to \mathfrak{g} \to \mathfrak{g} / \mathfrak{h}\) 平凡。
\end{enumerate}

断言 \(a\) 与 \(b\) 由命题 1 的刻画 (iii) 得出。部分 \(c\) 由 \(a\) 以及在约化李代数中每个理想都是直积因子这一事实（第I章§6第4节命题5的推论）得出。

命题 3. 令 G 为连通分支群有限的李群。下列条件等价：

\begin{enumerate}
\def\labelenumi{(\roman{enumi})}
\item
  李代数 \(\mathrm{L}(\mathrm{G})\) 紧。
\item
  群 \(\mathrm{Ad}(\mathrm{G})\) 紧。
\item
  L(G) 上存在在 G 的伴随表示下不变的分离的正对称双线性形式。
\end{enumerate}

*（iv）G 有在左平移与右平移下不变的黎曼度量。*

\begin{enumerate}
\def\labelenumi{(\roman{enumi})}
\item
  \(\Longrightarrow\) (ii)：若 L(G) 紧，则群 \(\mathrm{Ad}(\mathrm{G}_{0}) = \mathrm{Int}(\mathrm{L}(\mathrm{G}))\) 紧；由于它在 \(\mathrm{Ad}(\mathrm{G})\) 中的指标有限，后一群也紧。
\item
  \(\Longrightarrow\) (iii)：这由第1节得出。
\item
  \(\Longrightarrow\) (i)：由于 \(\operatorname{Int}(\mathrm{L}(\mathrm{G})) \subset \operatorname{Ad}(\mathrm{G})\)，这由命题 1 的刻画 (iii) 得出。
\end{enumerate}

\(^{*}\) (iii) \(\Longleftrightarrow\) (iv)：这由第III章§3第13节得出。*

\subsection{李代数为紧的李群}

定理 1.（H. Weyl）令 G 为李代数是紧半单李代数的连通李群。则 G 紧且其中心有限。

由于 G 半单，其中心 D 离散。此外，商群 G/D 同构于 \(\operatorname{Ad}(G)\)（第III章§6第4节推论4），因而是紧的（命题3）。最后，群 G/D 等于它的换位子群（第III章§9第2节命题4的推论）。于是定理由《积分学》第VII章§3第2节命题5 得出。

命题 4. 令 G 为李代数紧的连通李群。则存在一个环面 T、一个单连通紧半单李群 S、一个向量群 V，以及一个具有有限核的满射态射 \(f: V \times T \times S \to G\)。若 G 紧，则群 V 化为单位元。

令 C（分别地，S）为李代数同构于 L(G) 的中心（分别地，导代数）的单连通李群。则 C 是向量群，S 是中心有限的紧群（定理1），且 G 可与 C × S 关于一个离散子群 D 的商等同起来，D 是中心的（《积分学》第VII章§3第2节引理4）。由于 D 在 S 上投影的像中心因而有限，D∩C 在 D 中的指标有限。令 C' 为 C 中由 D ∩ C 生成的向量子空间，V 为其补子空间。则群 T = C'/(D ∩ C) 是环面，且 G 同构于积群 V × T × S 关于一个有限群的商。

若 G 紧，则 \(V \times T \times S\) 紧（《一般拓扑学》第III章§4第1节命题2的推论2），因而 V 紧，这就推出 \(V = \{e\}\)。

推论 1. 令 G 为连通紧李群。则 \(\mathrm{C}(\mathrm{G})_{0}\) 是环面，\(\mathrm{D}(\mathrm{G})\) 是连通紧半单李群，且态射 \((x,y)\mapsto xy\)（从 \(\mathrm{C}(\mathrm{G})_{0}\times\mathrm{D}(\mathrm{G})\) 到 G 的）是有限覆盖。

沿用命题4 的记号，我们有 \(\mathrm{V} = \{e\}\)，且子群 \(f(\mathrm{T})\) 与 \(f(\mathrm{S})\)（\(\mathrm{G}\) 的）紧，因而闭。于是只需证明 \(f(\mathrm{T}) = \mathrm{C}(\mathrm{G})_0\)、\(f(\mathrm{S}) = \mathrm{D}(\mathrm{G})\)。现在，\(\mathrm{L}(\mathrm{G}) = \mathrm{L}(f(\mathrm{T})) \times \mathrm{L}(f(\mathrm{S}))\)；由于 \(\mathrm{S}\) 半单而 \(\mathrm{T}\) 交换，这就推出 \(\mathrm{L}(f(\mathrm{T})) = \mathcal{C}(\mathrm{L}(\mathrm{G})) = \mathrm{L}(\mathrm{C}(\mathrm{G})_0)\)（第III章§9第3节命题8）且 \(\mathrm{L}(f(\mathrm{S})) = \mathcal{DL}(\mathrm{G}) = \mathrm{L}(\mathrm{D}(\mathrm{G}))\)（第III章§9第2节命题4的推论），由此即得所述断言。

推论 2. 连通紧半单李群的中心与基本群都是有限的。它的泛覆盖是紧的。

沿用命题4 的记号，群 V 与 T 都化为单位元；于是 S 是 G 的泛覆盖，且 G 的基本群同构于 \(\operatorname{Ker} f\)，因而有限。G 的中心 D 由于 G 半单而是离散的，故 D 有限。

推论 3. 连通紧李群 G 的基本群是有限型的 Z-模，其秩等于 C(G) 的维数。

事实上，沿用推论1 的记号，\(\mathrm{C}(\mathrm{G})_{0}\) 的基本群同构于 \(Z^{n}\)，其中 \(n = \dim \mathrm{C}(\mathrm{G})_{0}\)，而 \(\mathrm{D}(\mathrm{G})\) 的基本群有限（推论2）。

推论 4. 令 G 为连通紧李群。下列条件等价：

\begin{enumerate}
\def\labelenumi{(\roman{enumi})}
\item
  G 半单；
\item
  C(G) 有限；
\item
  \(\pi_{1}(G)\) 有限。
\end{enumerate}

若 \(\mathrm{G}\) 单连通，则它半单。

这由推论1 至 3 得出。

推论 5. 令 G 为连通紧李群。则 Int(G) 是李群 Aut(G) 的单位元连通分支（第III章§10第2节）。

令 \(f \in \mathrm{Aut}(\mathrm{G})_0\)。则 \(f\) 诱导出一个自同构 \(f_1\)（\(\mathrm{C}(\mathrm{G})_0\) 的）与一个自同构 \(f_2\)（\(\mathrm{D}(\mathrm{G})\) 的），且我们有 \(f_1 \in \mathrm{Aut}(\mathrm{C}(\mathrm{G})_0)_0\)，\(f_2 \in \mathrm{Aut}(\mathrm{D}(\mathrm{G}))_0\)。由于 \(\mathrm{Aut}(\mathrm{C}(\mathrm{G})_0)\) 离散（《一般拓扑学》第VII章§2第4节命题5），我们有 \(f_1 = \mathrm{Id}\)；由于 \(\mathrm{D}(\mathrm{G})\) 半单，由第III章§10第2节定理1的推论2，存在元素 \(g\)（\(\mathrm{D}(\mathrm{G})\) 的）使 \(f_2(x) = g x g^{-1}\) 对所有 \(x \in \mathrm{D}(\mathrm{G})\) 成立。对所有 \(x \in \mathrm{C}(\mathrm{G})_0\)，我们有 \(g x g^{-1} = x = f_1(x)\)；由于 \(\mathrm{G} = \mathrm{C}(\mathrm{G})_0.\mathrm{D}(\mathrm{G})\)，由此推出 \(g x g^{-1} = f(x)\) 对所有 \(x \in \mathrm{G}\) 成立，即 \(f = \operatorname{Int} g\)。

命题 5. 令 \(\mathbf{G}\) 为李代数紧的李群。

\begin{enumerate}
\def\labelenumi{\alph{enumi})}
\item
  假设 G 连通。则 G 有一个最大紧子群 K；它是连通的。存在 G 的一个闭的中心向量子群（第2节）N，使 G 是直积 \(N \times K\)。
\item
  假设 G 的连通分支群有限。则：
\end{enumerate}

\begin{enumerate}
\def\labelenumi{(\roman{enumi})}
\item
  \(\mathbf{G}\) 的每个紧子群都含于某个极大紧子群。
\item
  若 \(\mathrm{K}_1\) 与 \(\mathrm{K}_2\) 是 \(\mathbf{G}\) 的两个极大紧子群，则存在 \(g \in \mathbf{G}\) 使 \(\mathrm{K}_2 = g\mathrm{K}_1g^{-1}\)。
\item
  令 \(\mathrm{K}\) 为 \(\mathrm{G}\) 的极大紧子群。则 \(\mathrm{K} \cap \mathrm{G}_0\) 等于 \(\mathrm{K}_0\)；它是 \(\mathrm{G}_0\) 的最大紧子群。
\item
  存在 \(G_{0}\) 的一个闭的中心向量子群 N，它在 G 中正规，使得对 G 的任何极大紧子群 K，\(G_{0}\) 是 \(K_{0}\) 与 N 的直积，而 G 是 K 与 N 的半直积。
\end{enumerate}

\begin{enumerate}
\def\labelenumi{\alph{enumi})}
\item
  保留命题4 的记号。\(\mathrm{Ker}f\) 在 V 上的投影是向量群 V 的有限子群，故化为单位元。由此推出 \(\mathrm{Ker}f\) 含于 \(\mathrm{T} \times \mathrm{S}\)，从而 G 是向量群 \(\mathrm{N} = f(\mathrm{V})\) 与紧群 \(\mathrm{K} = f(\mathrm{T} \times \mathrm{S})\) 的直积。G 的每个紧子群在 N 上的投影都化为单位元，故含于 K。这就证明了 \(a\))。
\item
  现在假设 \(G/G_{0}\) 有限。由 a)，\(G_{0}\) 是其最大紧子群 M 与一个向量子群 P 的直积；G 的子群 M 显然正规。令 n 为 \(L(M)\) 在 \(L(G)\) 中的补向量子空间，它在 G 的伴随表示下稳定（第1节及第3节命题3）；它是 \(L(G)\) 的理想，且 \(L(G) = L(M) \times n\)。令 N 为 G 的以 n 为李代数的积分子群；由第III章§6第6节命题14，它在 G 中正规。\(L(G)\) 到 \(L(P)\) 上的以 \(L(M)\) 为核的投影诱导出 n 到 \(L(P)\) 的同构；由此推出 \(G_{0}\) 到 P 上的投影诱导出 N 到 P 的一个 étale 态射；由于 P 单连通，这是一个同构，故 N 是向量群。态射 \((x, y) \mapsto xy\)（从 \(M \times N\) 到 \(G_{0}\) 的）是单的 étale 态射（因为 \(M \cap N\) 化为单位元），因而是同构。由此推出 N 是 G 的闭子群，且商 G/N 紧，因为 \(G_{0}/N\) 紧而 \(G/G_{0}\) 有限（《一般拓扑学》第III章§4第1节命题2的推论2）。
\end{enumerate}

由《积分学》第VII章§3第2节命题3，G 的每个紧子群都含于某个极大紧子群，这些极大紧子群相互共轭，且对 G 的任何极大紧子群 K，G 是 K 与 N 的半直积。由于 \(G_{0}\) 含有 N，它是 N 与 \(G_{0} \cap K\) 的半直积；由此推出 \(G_{0} \cap K\) 连通，从而等于 \(K_{0}\)，因为 \(\mathrm{K}/(\mathrm{G}_{0} \cap \mathrm{K})\) 同构于 \(G/G_{0}\)，因而有限；最后，由 a)，\(K_{0}\) 显然是 \(G_{0}\) 的最大紧子群。

推论. 若 N 满足 b) (iv) 的条件，且 \(K_{1}\) 与 \(K_{2}\) 是 G 的两个极大紧子群，则存在 \(n \in N\) 使 \(nK_{1}n^{-1} = K_{2}\)。

事实上，由 (ii) 存在元素 \(g \in G\) 使 \(gK_{1}g^{-1} = K_{2}\)；由 (iv) 存在 \(n \in N\) 与 \(k \in K_{1}\) 使 g = nk。于是元素 n 具有所要的性质。

\section{紧李群的极大环面}